% =============================================================================
%  Incompletitud, formalizada
%  Julián Calderón Almendros · Licencia MIT (código) / CC BY-SA 4.0 (prosa)
%
%  Documento maestro. El índice sigue PLAN-LIBRO.md §3. Los capítulos aún no
%  escritos van comentados, para que el esqueleto refleje el plan completo.
%
%  ⚠️ Los ficheros de capítulo NO llevan número: el orden lo fija SÓLO este
%  fichero. Un `cap04-…` que acabara siendo el capítulo 8 es justo el tipo de
%  mentira silenciosa contra la que va todo el aparato de controles.
%
%  Compilar:  make          (extrae, compila y verifica)
% =============================================================================
\documentclass[11pt,openright]{book}
% =============================================================================
%  preamble.tex — Incompletitud, formalizada
%  Estilo, paquetes y los cinco entornos del libro.
%  Ver PLAN-LIBRO.md §4. NO contiene texto del libro.
% =============================================================================

% --- Motor -------------------------------------------------------------------
% El libro se compila con LuaLaTeX. Motivo: el proyecto usa Unicode
% intensivamente (σ ⟹ ⊢ ≐ ⌜·⌝ ⟨·⟩ Σ₁ Δ₀ ˙) y fontspec + luaotfload permiten
% una cadena de fuentes de respaldo por glifo. Con pdflatex esto no es viable.
\usepackage{iftex}
\ifLuaTeX\else\ifXeTeX\else
  \PackageError{preamble}{Compila con LuaLaTeX o XeLaTeX, no con pdfLaTeX}{%
    El libro necesita fontspec para el Unicode del proyecto.}%
\fi\fi

\usepackage[a4paper,margin=2.6cm,bottom=3cm]{geometry}
\usepackage{fontspec}
% babel español si la instalación de TeX lo trae; si no, se compila igual y
% sólo se pierden la partición de palabras y los rótulos automáticos.
\IfFileExists{spanish.ldf}{%
  \usepackage[spanish,es-noshorthands]{babel}%
}{%
  \PackageWarningNoLine{preamble}{spanish.ldf no disponible:
    instala texlive-lang-spanish para partición de palabras correcta}%
  \AtBeginDocument{%
    \renewcommand{\contentsname}{Índice}%
    \renewcommand{\partname}{Parte}%
    \renewcommand{\chaptername}{Capítulo}%
    \renewcommand{\appendixname}{Apéndice}%
    \renewcommand{\bibname}{Bibliografía}%
    \renewcommand{\tablename}{Tabla}%
    \renewcommand{\figurename}{Figura}%
  }%
}
\usepackage{csquotes}
\usepackage{microtype}
\usepackage{xcolor}
\usepackage{amsmath,amssymb,amsthm,mathtools}
\usepackage{booktabs}
\usepackage{tabularx}
\usepackage{array}
\usepackage{enumitem}
\usepackage{etoolbox}
\usepackage[most]{tcolorbox}
\usepackage{fancyvrb}
\usepackage{fvextra}  % breaklines para fancyvrb
\usepackage{graphicx}
\usepackage{hyperref}
\usepackage[nameinlink]{cleveref}

% Margen de estiramiento para párrafos con mucho \texttt: sin él, una racha de
% identificadores sin puntos de corte desborda por unos pocos puntos.
\setlength{\emergencystretch}{2.5em}

% --- Fuentes -----------------------------------------------------------------
% Texto: se prefiere Libertinus (buena cobertura griega y matemática);
% si no está instalada, Latin Modern.
\IfFontExistsTF{Libertinus Serif}{%
  \setmainfont{Libertinus Serif}\setsansfont{Libertinus Sans}%
}{\IfFontExistsTF{Latin Modern Roman}{%
  \setmainfont{Latin Modern Roman}\setsansfont{Latin Modern Sans}%
}{%
  \setmainfont{DejaVu Serif}\setsansfont{DejaVu Sans}%
}}

% Código: DejaVu Sans Mono cubre casi todo el Unicode del proyecto.
% Se declara una CADENA DE RESPALDO por glifo: si DejaVu no tiene el signo,
% luaotfload lo busca en las siguientes. Sin esto aparecen cuadraditos en
% ⌜·⌝, ≐ o los puntos de «dotar» (ȧ).
% Sólo LuaLaTeX admite cadena de respaldo por glifo (luaotfload). Bajo XeLaTeX
% se depende de la cobertura de la fuente elegida.
\makeatletter
\newif\ifrpp@fb  % ¿hay cadena de respaldo por glifo?
\ifLuaTeX
\IfFileExists{luaotfload-main.lua}{%
  \directlua{
    if luaotfload and luaotfload.add_fallback then
      luaotfload.add_fallback("leanfallback", {
        "DejaVuSansMono:mode=harf;",
        "Symbola:mode=harf;",
        "NotoSansMath-Regular:mode=harf;",
        "FreeMono:mode=harf;",
      })
    end
  }%
  \rpp@fbtrue
}{}
\fi
\IfFontExistsTF{DejaVu Sans Mono}{%
  \ifrpp@fb
    \setmonofont{DejaVu Sans Mono}[Scale=MatchLowercase,
                                   RawFeature={fallback=leanfallback}]%
  \else
    \setmonofont{DejaVu Sans Mono}[Scale=MatchLowercase]%
  \fi
}{%
  \setmonofont{Latin Modern Mono}[Scale=MatchLowercase]%
}
\makeatother

% --- Paleta ------------------------------------------------------------------
\definecolor{rppTinta}   {HTML}{1B1B1F}  % texto
\definecolor{rppCodigo}  {HTML}{F6F7F9}  % fondo de código
\definecolor{rppBorde}   {HTML}{C9CED6}
\definecolor{rppMate}    {HTML}{2F5D8C}  % matematica  — azul sobrio
\definecolor{rppLeccion} {HTML}{1F7A63}  % leccion     — verde azulado
\definecolor{rppMuro}    {HTML}{B0641E}  % muro        — ámbar: obstrucción
\definecolor{rppRefut}   {HTML}{9B2226}  % refutado    — rojo: resultado negativo
\definecolor{rppSondeo}  {HTML}{6E6E78}  % marca «fuera del build»
\definecolor{rppKw}      {HTML}{7A3E9D}
\definecolor{rppCom}     {HTML}{6E7781}
\definecolor{rppStr}     {HTML}{1F7A63}

\hypersetup{colorlinks=true, linkcolor=rppMate, citecolor=rppMate,
            urlcolor=rppMate, linktoc=all}

% --- Resaltado de Lean 4 -----------------------------------------------------
% NO se usa `listings`. Medido el 2026-09-03: con Unicode, listings REORDENA los
% caracteres —`(φ : Formula)` sale impreso como `φ( : Formula)`—, es decir,
% ALTERA EN SILENCIO el código publicado. En un libro cuyo principio editorial
% §2.1 es justo que eso no ocurra, es descalificante. Tampoco `minted`: exige
% --shell-escape y un lexer de Lean 4 en Pygments que no se puede dar por dado.
%
% Se usa `fancyvrb`, que es fiel byte a byte, y el RESALTADO LO HACE EL
% EXTRACTOR: scripts/extraer.py emite \kw{}/\cm{}/\st{} ya marcados. Es además
% lo correcto arquitectónicamente: quien conoce el código es el extractor.
\newcommand{\kw}[1]{\textcolor{rppKw}{\textbf{#1}}}   % palabra clave
\newcommand{\cm}[1]{\textcolor{rppCom}{\textit{#1}}}  % comentario
\newcommand{\st}[1]{\textcolor{rppStr}{#1}}           % cadena

% =============================================================================
%  MARCAS DE NIVEL — objeto vs meta, siempre a la vista
% =============================================================================
% Por qué esto no es decoración. La inconsistencia de ADR-012 fue un ERROR DE
% CATEGORÍA: `tcFn` es una operación sobre SINTAXIS declarada como función
% OBJETO. Si el nivel no se ve, el error no se ve. De ahí que cada fragmento de
% código y cada enunciado del libro lleve su nivel impreso.
%
% Y no es una distinción binaria: hay DOS EJES independientes.
%   · ¿QUIÉN afirma?  Lean (metateoría) · ⊢ (cálculo ω) · Prf (cálculo finitario)
%   · ¿SOBRE QUÉ?     Lean · sintaxis · objeto · código · Prov
% `prf_formCode_numeral` afirma en Prf SOBRE códigos; `hFN` afirma en ⊢ sobre lo
% mismo; `two_mul_consN` afirma en Lean sobre Nat. Son tres casillas distintas.

\definecolor{nvLean}  {HTML}{5B6470}   % metateoría: Lean
\definecolor{nvDer}   {HTML}{2F5D8C}   % el cálculo ⊢ (ω)
\definecolor{nvPrf}   {HTML}{1F7A63}   % el cálculo Prf (finitario, clásico)
\definecolor{nvPrf0}  {HTML}{3E9C84}   % Prf₀ (finitario, intuicionista)
\definecolor{nvPrfH}  {HTML}{15594A}   % PrfH (finitario CON contexto)
\definecolor{nvObj}   {HTML}{B0641E}   % lenguaje objeto
\definecolor{nvCod}   {HTML}{7A3E9D}   % códigos (términos objeto que denotan sintaxis)
\definecolor{nvProv}  {HTML}{9B2226}   % dentro de Prov

\makeatletter
\newcommand{\rpp@nvcolor}[1]{%
  \ifcsname nvcol@#1\endcsname \csname nvcol@#1\endcsname \else nvLean\fi}
\expandafter\def\csname nvcol@Lean\endcsname{nvLean}
\expandafter\def\csname nvcol@derives\endcsname{nvDer}
\expandafter\def\csname nvcol@Prf\endcsname{nvPrf}
\expandafter\def\csname nvcol@Prf0\endcsname{nvPrf0}
\expandafter\def\csname nvcol@PrfH\endcsname{nvPrfH}
\expandafter\def\csname nvcol@sintaxis\endcsname{nvLean}
\expandafter\def\csname nvcol@objeto\endcsname{nvObj}
\expandafter\def\csname nvcol@codigo\endcsname{nvCod}
\expandafter\def\csname nvcol@Prov\endcsname{nvProv}

\newcommand{\rpp@nvlabel}[1]{%
  \ifcsname nvlab@#1\endcsname \csname nvlab@#1\endcsname \else #1\fi}
\expandafter\def\csname nvlab@Lean\endcsname{Lean}
% La etiqueta lleva el TIPO a la vista: si el juicio arrastra contexto (Γ) o no.
% Confundir un juicio con contexto y uno sin él es un error de tipo, no de estilo.
\expandafter\def\csname nvlab@derives\endcsname{$\Gamma\vdash$}
\expandafter\def\csname nvlab@Prf\endcsname{Prf}
\expandafter\def\csname nvlab@Prf0\endcsname{Prf$_0$}
\expandafter\def\csname nvlab@PrfH\endcsname{$\Gamma\,$PrfH}
\expandafter\def\csname nvlab@sintaxis\endcsname{sintaxis}
\expandafter\def\csname nvlab@objeto\endcsname{objeto}
\expandafter\def\csname nvlab@codigo\endcsname{código}
\expandafter\def\csname nvlab@Prov\endcsname{Prov}

% \nivel{quién afirma}{sobre qué} — la chapa de dos casillas.
\newcommand{\nivel}[2]{%
  \begingroup\sffamily\scriptsize\bfseries
  \colorbox{\rpp@nvcolor{#1}}{\textcolor{white}{\strut\,\rpp@nvlabel{#1}\,}}%
  \colorbox{\rpp@nvcolor{#2}!18}{\textcolor{\rpp@nvcolor{#2}}{\strut\,\rpp@nvlabel{#2}\,}}%
  \endgroup}
\makeatother

% Marcas de nivel en prosa y en fórmulas.
\newcommand{\obj}[1]{\textcolor{nvObj}{#1}}        % símbolo del lenguaje objeto
\newcommand{\cod}[1]{\textcolor{nvCod}{#1}}        % código
\newcommand{\prov}[1]{\textcolor{nvProv}{#1}}      % dentro de Prov

% --- Nomenclatura (PLAN-LIBRO.md §2.7) --------------------------------------
% \defterm marca LA introducción de un término. El control `make terminos`
% exige que ningún término del vocabulario aparezca en el texto ANTES de su
% \defterm — sin exenciones: un lector con bases tiene que poder avanzar.
\newcommand{\defterm}[1]{\textbf{\textcolor{rppMate}{#1}}}

% Etiqueta de nivel de un TÉRMINO del glosario. Contesta, para cada palabra, la
% pregunta que el libro no puede dejar implícita: ¿esto es objeto o es meta?
\makeatletter
\newcommand{\nivelterm}[1]{%
  \begingroup\sffamily\scriptsize\bfseries
  \ifcsname ntcol@#1\endcsname
    \colorbox{\csname ntcol@#1\endcsname!16}{\textcolor{\csname ntcol@#1\endcsname}{\strut\,#1\,}}%
  \else
    \colorbox{nvLean!16}{\textcolor{nvLean}{\strut\,#1\,}}%
  \fi
  \endgroup}
\expandafter\def\csname ntcol@meta\endcsname{nvLean}
\expandafter\def\csname ntcol@objeto\endcsname{nvObj}
\expandafter\def\csname ntcol@meta→objeto\endcsname{nvCod}
\expandafter\def\csname ntcol@ambos\endcsname{nvDer}
\expandafter\def\csname ntcol@proyecto\endcsname{nvPrf}
\makeatother
\newcommand{\nota}[1]{{\footnotesize\color{rppCom} #1}}

% =============================================================================
%  LOS CINCO ENTORNOS
%  Ver PLAN-LIBRO.md §4. Cada uno tiene un papel distinto y no intercambiable.
% =============================================================================

% --- (1) lean — código Lean 4 real ------------------------------------------
% Uso normal: NO se escribe a mano. El extractor genera extraido/<clave>.tex y
% el capítulo lo trae con \leanfrag{clave} (§2.1: nada de copiar-pegar).
% El argumento opcional es la referencia de origen (módulo:declaración).
\newtcolorbox{leanbox}[1][]{%
  enhanced, breakable, sharp corners,
  colback=rppCodigo, colframe=rppBorde, boxrule=0.4pt,
  left=2mm, right=2mm, top=1.4mm, bottom=1.4mm,
  fonttitle=\ttfamily\scriptsize, coltitle=rppTinta,
  attach boxed title to top left={xshift=2mm, yshift=-2mm},
  boxed title style={colback=rppBorde, sharp corners, boxrule=0pt},
  #1}

\DefineVerbatimEnvironment{lean}{Verbatim}{%
  fontsize=\small, commandchars=\\\{\},
  breaklines=true, breakanywhere=false, breakautoindent=true,
  % El símbolo de continuación NO puede ser una flecha: pdftotext la extrae como
  % `→`, que es un token legítimo de Lean — un lector podría leer una flecha de
  % más en una línea partida, y el verificador del PDF no podría distinguirlas.
  % Se usa «›», que no aparece ni una vez en el código del proyecto ni en FOL.
  breaksymbolleft=\textcolor{rppBorde}{\guilsinglright},
  breaksymbolright={},
  xleftmargin=0pt, samepage=false}

% Trae un fragmento ya extraído del repo.
\newcommand{\leanfrag}[1]{\input{extraido/#1.tex}}

% Referencia de origen impresa bajo un fragmento.
\newcommand{\leanorigen}[2]{%
  \par\vspace{-0.4\baselineskip}%
  {\ttfamily\scriptsize\color{rppCom} #1\ \textbf{·}\ #2\par}}

% Footprint de `#print axioms` — OBLIGATORIO en todo teorema (§2.2(c)).
\newcommand{\footprint}[1]{%
  \par\vspace{-0.2\baselineskip}%
  {\ttfamily\scriptsize\color{rppCom}\#print axioms $\rightarrow$ #1\par}}

% Marca «fuera del build» para fragmentos de sondeos/ (§2.3).
\newcommand{\fueradelbuild}[1]{%
  \par\vspace{-0.2\baselineskip}%
  {\sffamily\scriptsize\color{rppSondeo}%
   \fcolorbox{rppSondeo}{white}{\strut\,fuera del build\,}\quad
   recompilar con \texttt{lake env lean #1}\par}}

% --- (2) matematica — el enunciado formal, legible sin Lean ------------------
% Sin barra de título: el rótulo va como línea de texto, para que el enunciado
% se lea como un enunciado y no como un aviso.
\NewDocumentEnvironment{matematica}{O{}}{%
  \begin{tcolorbox}[enhanced, breakable, sharp corners,
    colback=white, colframe=white, boxrule=0pt,
    borderline west={1.6pt}{0pt}{rppMate},
    left=4mm, right=2mm, top=1.6mm, bottom=1.6mm]%
  \ifblank{#1}{}{{\sffamily\bfseries\small\color{rppMate} #1\par}\vspace{0.4em}}%
}{%
  \end{tcolorbox}%
}

% --- (3) leccion — la nota pedagógica ---------------------------------------
\newtcolorbox{leccion}[1][Lección]{%
  enhanced, breakable, sharp corners,
  colback=rppLeccion!4, colframe=rppLeccion, boxrule=0.5pt,
  left=3mm, right=3mm, top=1.8mm, bottom=1.8mm,
  fonttitle=\bfseries\sffamily\small, coltitle=white,
  colbacktitle=rppLeccion,
  attach boxed title to top left={xshift=3mm, yshift=-2.4mm},
  boxed title style={sharp corners, boxrule=0pt},
  title={#1}}

% --- (4) muro — una obstrucción encontrada ----------------------------------
\newtcolorbox{muro}[1][Muro]{%
  enhanced, breakable, sharp corners,
  colback=rppMuro!5, colframe=rppMuro, boxrule=0.5pt,
  left=3mm, right=3mm, top=1.8mm, bottom=1.8mm,
  fonttitle=\bfseries\sffamily\small, coltitle=white,
  colbacktitle=rppMuro,
  attach boxed title to top left={xshift=3mm, yshift=-2.4mm},
  boxed title style={sharp corners, boxrule=0pt},
  title={#1}}

% --- (5) refutado — un resultado NEGATIVO, compilado ------------------------
% No es un muro: un muro es «no sé pasar»; esto es «se ha probado que no se
% pasa». Es el entorno que sostiene el capítulo 25 y media Parte IV.
\newtcolorbox{refutado}[1][Refutado]{%
  enhanced, breakable, sharp corners,
  colback=rppRefut!5, colframe=rppRefut, boxrule=0.5pt,
  left=3mm, right=3mm, top=1.8mm, bottom=1.8mm,
  fonttitle=\bfseries\sffamily\small, coltitle=white,
  colbacktitle=rppRefut,
  attach boxed title to top left={xshift=3mm, yshift=-2.4mm},
  boxed title style={sharp corners, boxrule=0pt},
  title={#1}}

% =============================================================================
%  Entornos de teorema (notación matemática estándar)
% =============================================================================
\theoremstyle{plain}
\newtheorem{teorema}{Teorema}[chapter]
\newtheorem{lema}[teorema]{Lema}
\newtheorem{proposicion}[teorema]{Proposición}
\newtheorem{corolario}[teorema]{Corolario}
\theoremstyle{definition}
\newtheorem{definicion}[teorema]{Definición}
\newtheorem{convenio}[teorema]{Convenio}
\theoremstyle{remark}
\newtheorem*{observacion}{Observación}

% =============================================================================
%  Macros de prosa
% =============================================================================
% \detokenize: permite escribir \ident{ax_tc_cons} sin escapar el subrayado.
% (El nombre \lean está tomado por el ENTORNO `lean`; en prosa se usa \ident.)
\newcommand{\ident}[1]{\texttt{\detokenize{#1}}}         % identificador Lean
% Las rutas de módulo son largas y aparecen constantemente. Con \texttt son
% inquebrables y desbordan la caja de texto; \nolinkurl parte por «/» y «.»,
% que es donde una ruta se puede partir sin confundir al lector.
\newcommand{\modulo}[1]{\begingroup\urlstyle{tt}\nolinkurl{#1}\endgroup}
\newcommand{\adr}[1]{\textsc{adr}-#1}                    % referencia a un ADR
\newcommand{\printaxioms}{\texttt{\#print axioms}}       % \detokenize duplica #
\newcommand{\doc}[1]{\textsf{\detokenize{#1}}}           % fichero .md del repo
\newcommand{\codigo}[1]{\ensuremath{\ulcorner #1 \urcorner}}  % ⌜φ⌝
\newcommand{\dotado}[1]{\ensuremath{\dot{#1}}}           % ȧ  («dotar»)
\newcommand{\Prov}{\operatorname{Prov}}
\newcommand{\Prf}{\mathsf{Prf}}
\newcommand{\derives}{\vdash}
\newcommand{\noderives}{\nvdash}

% Aviso editorial recurrente (§6): se usa donde toque, no se improvisa.
\newcommand{\avisoreparacion}{%
  \begin{muro}[Aviso]
  Lo que se retiró fue la inconsistencia \textbf{conocida y localizada}
  (\ident{ax_tc_cons}). \textbf{Eso no es una prueba de consistencia} de
  Q\texttt{++} extendido.
  \end{muro}}
\newcommand{\avisomediagodel}{%
  \begin{muro}[Aviso]
  Gödel I está \textbf{a medias}: sólo $\noderives G$. La indecidibilidad
  ($\noderives \lnot G$) \textbf{no está cerrada} en la cadena real.
  \end{muro}}


\title{\bfseries Incompletitud, formalizada\\[0.4em]
       \large Un recorrido por la formalización real de los teoremas de Gödel
       en Lean 4}
\author{Julián Calderón Almendros}
\date{\normalsize\vspace{2em}
  Prosa bajo \textbf{CC BY-SA 4.0} \quad·\quad código Lean 4 bajo \textbf{MIT}\\[0.3em]
  {\small Ver el apéndice \emph{Licencia}}}

\begin{document}

\frontmatter
\maketitle

% -- Estado del código sobre el que se escribe este ejemplar ------------------
% GENERADO por scripts/extraer.py. Es el control §2.2(b): el libro declara
% contra qué commit y con qué build se ha verificado cada teorema que cita.
\IfFileExists{extraido/_estado.tex}{\input{extraido/_estado.tex}}{%
  \chapter*{Estado del código}
  \textbf{(sin generar)} — ejecuta \texttt{make extraer} antes de compilar.}

\tableofcontents

\mainmatter

% =============================================================================
\part{El problema}
\label{parte:problema}

\noindent \emph{¿De dónde sale la pregunta, y qué hay que construir para contestarla con una máquina?} Esta parte no supone nada y no demuestra nada: sitúa el problema, distingue los lenguajes que el libro va a manejar y da las cuatro palabras de Lean que hacen falta para leer el resto.
\chapter{1931: qué preguntó Hilbert y qué contestó Gödel}
\label{cap:hilbert}

Este libro cuenta una formalización, pero la formalización no se entiende sin la pregunta que la
originó. Y la pregunta no la hizo Gödel: la hizo Hilbert, treinta años antes, y era una pregunta
optimista.

\section{El programa de Hilbert}

A principios del siglo~XX las matemáticas habían pasado un susto. Las paradojas de la teoría de
conjuntos mostraron que razonar con confianza sobre objetos infinitos podía llevar a
contradicciones, y que la intuición no bastaba para detectarlas. La respuesta de Hilbert fue un
programa: si no podemos fiarnos de la intuición, escribamos las matemáticas como un sistema formal
—símbolos, axiomas, reglas— y demostremos \emph{sobre ese sistema}, con métodos que nadie discuta,
que es consistente.

\begin{leccion}[La palabra clave es «finitario»]
Hilbert no pedía cualquier demostración de consistencia. Pedía una \textbf{finitaria}: que no usara
ella misma el infinito que estaba tratando de justificar. Razonar sobre secuencias finitas de
símbolos, comprobar cosas caso a caso, contar. Nada más. Ésa era la condición, y es la que da
sentido a todo lo que viene después — incluida buena parte de las decisiones de este proyecto.
\end{leccion}

Para plantear siquiera la pregunta hay que separar dos cosas que en matemáticas ordinarias van
juntas. Por un lado, el sistema del que se habla: sus fórmulas, sus axiomas, su noción de
demostración. Lo llamaremos la \defterm{teoría objeto}, y sus fórmulas son, para nosotros,
\emph{datos}: cadenas de símbolos que se pueden manipular sin creerlas. Por otro, el lugar desde el
que hablamos de ella y demostramos cosas \emph{sobre} ella: la \defterm{metateoría}.

Esa separación no es un tecnicismo. Es la condición de posibilidad del programa entero: sin ella,
«demostrar que la aritmética es consistente» sería demostrarlo \emph{dentro} de la aritmética, y no
convencería a nadie.

\section{El finitismo era una estrategia, no una creencia}
\label{sec:finitismo-estrategia}

Conviene deshacer ahora un malentendido que se arrastra, porque cambia cómo se lee todo lo demás.
Hilbert \textbf{no era} un matemático constructivista que desconfiara del infinito. Era justo lo
contrario: había hecho carrera usando métodos fuertemente no constructivos y los defendía como un
territorio ganado, un derecho que las matemáticas no debían devolver.

El caso famoso es el teorema de la base, en teoría de invariantes: Hilbert demostró que existe una
base finita \emph{sin exhibir ninguna}, en un campo cuyo rey indiscutible, Paul Gordan, había
construido las suyas a mano. La reacción que se cita siempre es la de Gordan —«esto no es
matemática, esto es teología»—, y este libro, que se dedica a comprobar afirmaciones, tiene que
añadir en el mismo aliento que \textbf{esa cita está mal atestiguada}: aparece impresa por primera
vez en la necrológica que Max Noether escribió \emph{veinticinco años después} del episodio. Lo que
sí está documentado es lo contrario del mito: Gordan publicó en 1893 su propia versión del
argumento diciendo que la demostración de Hilbert era, en sustancia, correcta.

\begin{observacion}
Se cita esa frase para ilustrar la resistencia al método no constructivo, y funciona tan bien como
ilustración que casi nadie va a mirar de dónde sale. Es exactamente el modo de fallo que
\doc{PLAN-LIBRO.md}~§2.6 describe para los docstrings de este proyecto: una afirmación cómoda,
repetida hasta ser evidente, cuya fuente nadie ha vuelto a abrir. Se cita aquí porque es
instructiva, y con la salvedad delante.
\end{observacion}

Con eso en su sitio, la exigencia finitaria se lee bien. Hilbert no pedía que las
\emph{matemáticas} fueran finitarias — pedía que lo fuera la \textbf{demostración de
consistencia}, que es otra cosa. Se puede usar todo el infinito que se quiera \emph{dentro} del
sistema formal, con tal de que el argumento que certifica ese sistema desde fuera no lo use. Es una
autolimitación de la metateoría, y es estratégica: sólo un argumento que no presuponga lo que está
en duda puede convencer a quien lo duda.

\begin{leccion}[Y Gödel hizo lo mismo, por la misma razón]
La restricción de Gödel también era estratégica. Su argumento no usa el infinito ni necesita lógica
clásica más allá de lo imprescindible — no porque él dudara de ellos, sino porque un teorema que
demoliera el programa de Hilbert usando los métodos que el programa quería justificar no habría
demolido nada. Los dos se autolimitaron para poder hablarle al otro.

Esa es, tres capas más abajo, la misma razón por la que este proyecto anota debajo de cada teorema
de qué axiomas depende, y por la que el capítulo~\ref{cap:objeto-y-meta} se molesta en contar
cuánta lógica clásica gasta:
\textbf{importa con qué se demuestra, no sólo qué se demuestra.}
\end{leccion}

\section{La idea de Gödel}

Gödel publicó en 1930 un resultado que encajaba con el optimismo: la lógica de primer orden es
completa —todo lo que es verdad en todos los modelos es demostrable—. Un año después publicó el
que no encajaba.

La idea central es una jugada de una audacia difícil de exagerar. Si las fórmulas de la teoría
objeto son cadenas de símbolos, y las cadenas de símbolos se pueden numerar, entonces \textbf{la
teoría puede hablar de sus propias fórmulas} — no de fórmulas, sino de los números que las
codifican. A ese procedimiento se le llama \defterm{aritmetización} o \defterm{gödelización}, y el
número —o el término que lo denota— es el \defterm{código de Gödel} de la fórmula.

\begin{leccion}[Por qué eso lo cambia todo]
Una teoría aritmética habla de números. Si las fórmulas \emph{son} números, la teoría habla de
fórmulas sin saberlo. Y si además puede expresar «existe una demostración de la fórmula número
$x$», entonces habla de su propia demostrabilidad. La autorreferencia deja de ser un juego de
palabras y pasa a ser aritmética.
\end{leccion}

Con esa maquinaria, Gödel construye una sentencia $G$ que la teoría demuestra equivalente a «yo no
soy demostrable» — un \defterm{punto fijo} de la negación de la demostrabilidad — y prueba que si la
teoría es consistente, no puede demostrar $G$. El primer teorema. Y a continuación observa que el
argumento entero, siendo finitario, se puede formalizar dentro de la propia teoría: de donde
resulta que la teoría tampoco puede demostrar su propia consistencia. El segundo teorema.

\section{El «I» del título}

Y aquí está el detalle que da sentido a este libro.

El artículo de 1931 se titula \emph{Sobre proposiciones formalmente indecidibles de los Principia
Mathematica y sistemas afines}, \textbf{I}. Ese «I» no es decorativo: Gödel anunció una segunda
parte con la demostración detallada del segundo teorema, que en el artículo publicado sólo queda
esbozada. \textbf{La segunda parte nunca se escribió.}

Lo que faltaba no era la idea, que estaba clara, sino la comprobación penosa de que el predicado de
demostrabilidad se comporta como debe: que si algo es demostrable, la teoría demuestra que lo es;
que la demostrabilidad respeta el modus ponens; y que la teoría demuestra ella misma esa primera
propiedad. Esas tres son las \defterm{condiciones de derivabilidad}, y las suministraron
Hilbert y Bernays en 1939.

\begin{leccion}[La tesis de este libro, en una frase]
Las condiciones de derivabilidad son el artículo que Gödel no llegó a publicar. Y de las tres, la
tercera —la que dice que la teoría demuestra \emph{de sí misma} que reconoce lo que es cierto de
forma finita, la \defterm{$\Sigma_1$-completitud} en su versión \defterm{provable}— es la que este
proyecto lleva meses sin poder cerrar. Buena parte de la Parte~\ref{parte:noescrito} cuenta ese intento.
\end{leccion}

\section{Qué quedó del programa}

El programa de Hilbert, tal y como estaba formulado, no se puede completar. Pero decir sólo eso es
quedarse corto de una manera que este libro quiere evitar: lo que Gödel demostró no es que las
matemáticas sean poco fiables, sino que \textbf{la fiabilidad no se puede certificar desde dentro}.
Es un resultado sobre los límites de la autocertificación, no sobre la verdad matemática.

Y tiene una consecuencia práctica que hoy se ve mejor que en 1931: si un sistema no puede
certificarse a sí mismo, la alternativa es certificarlo desde otro sistema — explícito, más débil o
más fuerte, pero declarado. Que es exactamente lo que hace una formalización asistida por ordenador,
y lo que hace este proyecto.

\chapter{Qué hay que construir para decirlo hoy}
\label{cap:construir}

El argumento del capítulo anterior cabe en diez páginas de prosa. Formalizarlo —hacer que una
máquina lo verifique paso a paso— exige construir cinco cosas, ninguna de las cuales se puede
saltar. Este capítulo dice cuáles son y en qué orden vienen. Es, de paso, el mapa del libro.

\section{Las cinco piezas}

\begin{enumerate}[leftmargin=1.6em]
\item \textbf{Un lenguaje.} Qué es un término, qué es una fórmula, cómo se sustituye una variable.
  Suena trivial y no lo es: la sustitución con nombres exige evitar capturas, y ahí se pierden más
  demostraciones de las que uno espera.
\item \textbf{Una teoría.} Un lenguaje no afirma nada. Hacen falta axiomas — y decidir cuáles, que
  es una decisión con consecuencias que este libro persigue hasta el final.
\item \textbf{Un cálculo.} Qué cuenta como demostración. Y aquí aparece la primera sorpresa del
  proyecto: hacen falta \emph{dos}, por una razón que el capítulo siguiente explica.
\item \textbf{Una codificación.} Convertir fórmulas en números, de modo que la teoría pueda hablar
  de ellas.
\item \textbf{Un verificador.} Una fórmula del propio lenguaje que diga «esto de aquí es una
  demostración correcta». Sin ella no hay predicado de demostrabilidad, y sin predicado de
  demostrabilidad no hay teorema de Gödel.
\end{enumerate}

\begin{leccion}[El orden no es negociable]
Cada pieza usa las anteriores. No se puede codificar sin lenguaje, ni verificar sin codificación.
Por eso el libro va en ese orden, y por eso los capítulos técnicos empiezan diciendo qué se sabe ya
— si el lector llegó hasta aquí, tiene derecho a que se le diga dónde está.
\end{leccion}

\section{Lo que este libro añade a esa lista}

Si el libro fuera sólo la construcción, sería un manual más. Hay dos cosas que lo separan de un
manual, y conviene anunciarlas ahora porque condicionan cómo está escrito.

\textbf{La primera: aquí todo está compilado.} Ningún bloque de código de este libro se ha escrito a
mano. Se extraen del repositorio, y hay controles que fallan si el texto impreso deja de coincidir
con el código, si un teorema deja de estar demostrado desde la base declarada, o si el libro nombra
un símbolo que no existe. Debajo de cada enunciado hay una línea gris con los axiomas de los que
depende — el capítulo~\ref{cap:lean} explica qué es eso y por qué está ahí.

\textbf{La segunda: aquí también se cuentan los fracasos.} Una formalización real produce, además de
teoremas, una cantidad notable de callejones sin salida — y algunos de ellos enseñan más que los
teoremas. La Parte~\ref{parte:noescrito} está dedicada a eso: una inconsistencia que sobrevivió meses, cuatro
reparaciones que no funcionan, y la que sí. En los manuales eso no aparece, porque los manuales
presentan el resultado ya limpio.

\begin{leccion}[Por qué los fracasos son contenido y no anécdota]
Un resultado negativo demostrado —\emph{esto no se puede hacer, y aquí está la prueba}— ahorra al
siguiente el tiempo de intentarlo. En este proyecto hay varios, compilados. El libro los trata con
el mismo rango que los positivos, y tiene un entorno tipográfico propio para ellos.
\end{leccion}

\section{Cómo leer este libro}

Las partes de este libro siguen las cinco piezas, más los teoremas y más lo que no sale en los
libros:

\begin{center}\small
\begin{tabularx}{\linewidth}{@{}l >{\raggedright\arraybackslash}X@{}}
\toprule
\textbf{I} · El problema & dónde estás ahora: el contexto y las cuatro palabras de Lean que hacen falta \\
\textbf{II} · El terreno & el lenguaje, los cálculos y la teoría aritmética \\
\textbf{III} · Aritmetización & la codificación y el verificador \\
\textbf{IV} · Los teoremas & el punto fijo, Gödel~I, las condiciones de derivabilidad, Gödel~II \\
\textbf{V} · Lo que no sale en los libros & la inconsistencia, las reparaciones, el método \\
\bottomrule
\end{tabularx}
\end{center}

\noindent Los apéndices recogen el inventario de axiomas, el mapa de módulos, las trampas de Lean
encontradas, y un glosario y una tabla de notación a los que se puede volver en cualquier momento.

\chapter{Los lenguajes de este libro}
\label{cap:lenguajes}

Hay una pregunta que en este libro se plantea en cada línea: \emph{esto que estoy leyendo,
¿se afirma dentro del sistema del que hablamos, o desde fuera?} Contestarla mal es caro — la Parte~\ref{parte:noescrito} cuenta una inconsistencia que fue exactamente
eso, un error de categoría. Este capítulo la contesta de una vez, y con más detalle del que suele
darse: porque ni hay una metateoría, ni hay un solo lenguaje del que se hable. Hay tres y hay
cuatro.

\section{Tres metateorías, no una}

Cuando se dice «la metateoría» se está diciendo, en tres momentos distintos de esta historia, tres
cosas incompatibles.

\textbf{La de Gödel era deliberadamente débil.} Razonamiento finitario sobre secuencias de símbolos,
y nada más — porque el objetivo era convencer a Hilbert, y una demostración de consistencia que
usara el infinito no habría valido de nada. Esa autolimitación no era humildad ni convicción
—§\ref{sec:finitismo-estrategia} cuenta que ninguno de los dos era constructivista—: era el
requisito para que el argumento le sirviera a quien iba dirigido.

\textbf{La del kernel de este proyecto es más fuerte de lo que parece.} El cálculo $\vdash$ de
\modulo{FOL} incluye la \defterm{$\omega$-regla}: de «$A[n]$ para todo término $n$» concluye
$\forall A$. Es válida en el modelo estándar y tiene \emph{infinitas premisas}, luego el cálculo que
la admite no es \defterm{r.e.} — no hay programa que liste sus teoremas.

\begin{muro}[La consecuencia que gobierna el libro entero]
Con la $\omega$-regla, $\vdash$ no es recursivamente enumerable. Y por el teorema de Tarski, ninguna
fórmula puede representar fielmente la demostrabilidad de un cálculo así. \textbf{Por eso el
proyecto construye un segundo cálculo}, finitario, y aritmetiza ése. Todo lo que en este libro se
codifica, se verifica o se diagonaliza ocurre en el segundo, no en el primero.
\end{muro}

\textbf{Y la de Lean es enormemente más fuerte que ambas.} Su teoría de tipos interpreta a ZFC y
más. Ahora bien —y esto sí está medido sobre el árbol—, \textbf{el proyecto usa dos niveles de ella
y ninguno más}: proposiciones y datos. No hay universos por encima, no hay polimorfismo de
universos, y no se cuantifica sobre proposiciones arbitrarias en ninguna parte. La fuerza está
disponible; no se gasta.

\subsection*{Dónde cae cada una: la escala, con nombres}

«Más fuerte» y «más débil» dicho así no informa. Hay una escala estándar, y las tres piezas de este
libro tienen sitio en ella. De arriba abajo por lo que \emph{demuestran} de aritmética:

\begin{center}\small
\begin{tabularx}{\linewidth}{@{}l >{\raggedright\arraybackslash}X l@{}}
\toprule
sistema & qué permite & papel aquí \\
\midrule
PRA & recursión primitiva, inducción sin cuantificadores no acotados & el finitismo que pedía Hilbert \\
$\mathrm{I}\Sigma_1$ & inducción restringida a fórmulas $\Sigma_1$ & lo que \emph{basta} como metateoría (cap.~\ref{cap:objeto-y-meta}) \\
HA & PA con lógica intuicionista & equiconsistente con PA \\
PA & inducción sobre todas las fórmulas de primer orden & la referencia habitual \\
Q\texttt{++} $+\ \omega$-regla & el cálculo $\vdash$ de este proyecto & \textbf{toda la aritmética verdadera} \\
\bottomrule
\end{tabularx}
\end{center}

La última fila no es una exageración retórica, y merece decirse con precisión porque es la razón
de ser del segundo cálculo.

\begin{matematica}[La $\omega$-regla no sube un escalón: se sale de la escala]
Es un resultado clásico que $\mathrm{PA}^{\omega}$ —PA cerrada bajo la $\omega$-regla— demuestra
\emph{exactamente} las sentencias verdaderas en $\mathbb{N}$:
\[ \mathrm{PA}^{\omega} \;=\; \mathrm{Th}(\mathbb{N}). \]
Y la demostración no usa casi nada de PA: le basta con que los axiomas sean verdaderos en
$\mathbb{N}$ y que las sentencias atómicas cerradas estén decididas. \textbf{Vale igual partiendo
de la aritmética de Robinson}, que es de donde parte este proyecto.
\end{matematica}

\selee{«añadir la $\omega$-regla a una teoría aritmética verdadera la convierte en la teoría de
\emph{todo} lo que es cierto de los números»}

O sea: $\vdash$ no está «un poco por encima de PA». Está en el tope, y por eso no es
\defterm{r.e.}: no hay nada que enumerar cuando ya se demuestra todo lo verdadero. El precio de esa
potencia es que sus demostraciones no son objetos finitos, y por tanto \textbf{no se pueden
codificar} — que es justo lo que la incompletitud necesita hacer.

\begin{muro}[La salvedad que hay que declarar]
El enunciado de arriba es sobre el modelo estándar. El vocabulario de este proyecto no es sólo el
de la aritmética: incluye los 107 axiomas de la capa de codificación, con símbolos opacos cuyo
significado fija la propia teoría. La igualdad $\mathrm{Th}(\mathbb{N})$ se lee, entonces, sobre
$\mathbb{N}$ con esos símbolos interpretados como se pretende. No cambia la moraleja —$\vdash$ no
es r.e., y por eso hace falta el segundo cálculo— pero sí la letra pequeña.
\end{muro}

\section{Cuatro lenguajes objeto, no uno}

Del otro lado pasa algo parecido, y es más fácil de pasar por alto porque los cuatro se llaman
igual. Un \defterm{lenguaje objeto} es sólo una gramática: dice cómo se escriben términos y
fórmulas, y no afirma nada. La teoría aparece cuando se le añaden axiomas.

\begin{center}\small
\begin{tabularx}{\linewidth}{@{}l >{\raggedright\arraybackslash}X >{\raggedright\arraybackslash}X r@{}}
\toprule
estrato & qué añade & para qué sirve & axiomas \\
\midrule
FOL$^=$ desnudo & la gramática y la lógica & escribir fórmulas; no afirma nada & --- \\
Q\texttt{++} & la \defterm{signatura} aritmética y sus \defterm{axiomas objeto} & \textbf{ser incompleta}: es la teoría de la que habla el teorema & 34 \\
Q\texttt{++} codificante & los símbolos opacos del verificador y sus ecuaciones & \textbf{hablar de sintaxis}: sin esto no hay predicado de demostrabilidad & +107 \\
con inducción (\modulo{Full/}) & el esquema de inducción & demostrar D2 y D3 & +1 esquema \\
\bottomrule
\end{tabularx}
\end{center}

\subsection*{Los dos del medio no se parecen en nada, y la diferencia decide el teorema}

Los estratos primero y último se distinguen solos. Los dos de en medio se confunden —los dos se
llaman «Q\texttt{++}», los dos son listas de fórmulas postuladas— y son de naturaleza distinta.

Los \textbf{34} son aritmética: hablan de sumar, multiplicar, comparar. Se entienden mirándolos, se
discuten como se discute un axioma, y son los que hacen del sistema \emph{una teoría de números},
que es lo que el teorema de Gödel exige para tener algo que decir. El
capítulo~\ref{cap:qpp} los desmenuza uno a uno.

Los \textbf{107} no hablan de números: hablan de \emph{sintaxis codificada como} números. Son las
ecuaciones que fijan el comportamiento de símbolos que no se pueden definir —\ident{substfc},
\ident{lenc}, \ident{nthc}, \ident{carc}, el verificador entero— y su papel no es ser verdad sobre
$\mathbb{N}$, sino \textbf{dar a la teoría un vocabulario para hablar de sus propias fórmulas}. Sin
ellos, el predicado «$x$ es demostrable» no se puede ni escribir dentro del lenguaje. El
capítulo~\ref{cap:verificador} los construye.

\begin{observacion}
La proporción es el dato: \textbf{tres veces más maquinaria que aritmética}. Eso es, medido, lo que
cuesta que una teoría pueda hablar de su propia sintaxis — y es la cifra que más sorprende a quien
llega desde la exposición clásica, donde la codificación se despacha con un «y esto se aritmetiza
de manera rutinaria».
\end{observacion}

\begin{muro}[Y por qué la frontera entre los dos no es cosmética]
El verificador consulta la lista de axiomas de la teoría: la regla que acepta «esto es un axioma»
comprueba que el código de la fórmula está en esa lista. Luego \textbf{la lista entra en la
definición del predicado de demostrabilidad}, y ese predicado entra en la sentencia de Gödel. Mover
un axioma de fuera a dentro no hace la teoría «un poco más fuerte»: cambia el predicado, y con él
cambia $G$. Es literalmente otro teorema.

Por eso este libro dice siempre de qué estrato habla, y por eso \adr{015} decidió \emph{definir}
la buena formación de los códigos en vocabulario ya existente en vez de postularla. La historia
completa está en \capfuturo{definir en vez de axiomatizar}.
\end{muro}

\begin{leccion}[Cuatro teorías, cuatro teoremas de incompletitud]
No son cuatro presentaciones de lo mismo: son cuatro teorías distintas, y el teorema de
incompletitud que se demuestra sobre cada una es un teorema distinto. En cuanto la teoría contiene
un verificador que consulta su propia lista de axiomas, \textbf{cambiar la lista cambia la sentencia
de Gödel}. El libro vuelve sobre esto dos veces, y las dos con consecuencias prácticas.
\end{leccion}

\section{Los signos que hacen falta desde ya}

Cuatro notaciones aparecen a partir de aquí constantemente, y conviene fijarlas antes.

\begin{center}\small
\begin{tabularx}{\linewidth}{@{}l >{\raggedright\arraybackslash}X@{}}
\toprule
$\Gamma \vdash A$ & $A$ se deriva de $\Gamma$ en el cálculo $\omega$ \\
$\defnot{\Prf}\,A$ & $A$ se deriva en el cálculo de Hilbert finitario, sin hipótesis
  \selee{«$A$ tiene una demostración finitaria»} \\
$\defnot{\codigo{\varphi}}$ & el \defterm{código de Gödel} de $\varphi$, \textbf{calculado desde
  fuera}: se lo damos ya hecho a la teoría \selee{«el código de $\varphi$»} \\
$\defnot{\dotado{a}}$ & «$a$ dotado»: el código de $a$, \textbf{calculado dentro} por una función
  del propio lenguaje \selee{«$a$ dotado», o «el código de lo que valga $a$»} \\
$\defnot{\Prov}(\codigo{\varphi})$ & la fórmula \emph{objeto} que dice «$\varphi$ es demostrable»
  \selee{«es demostrable», dicho dentro de la teoría} \\
\bottomrule
\end{tabularx}
\end{center}

\noindent Y dos convenios del lenguaje que sorprenden la primera vez. Las variables se escriben
como \defterm{índice de De Bruijn} —\texttt{\#0} es la ligada más cercana, y los cuantificadores no
llevan nombre de variable—. Y el término que denota el número $n$ es el \defterm{numeral}
$\sigma^n 0$: es \emph{canónico}, sólo hay una forma de escribirlo, y de ahí que sirva para
representar códigos.

\subsection*{Dos maneras de decir «el código de», y no son intercambiables}

Las dos filas segunda y tercera parecen decir lo mismo y no lo dicen. La diferencia es \emph{quién
hace el cálculo}, y de ella depende medio libro.

$\codigo{\varphi}$ es una \textbf{función de Lean}: nosotros, desde fuera, tomamos una fórmula
concreta y calculamos su código. El resultado es un término cerrado y concreto que se puede escribir
en un papel. Para usarlo hay que tener la fórmula \textbf{en la mano}.

$\dotado{a}$ es un \textbf{símbolo de función del lenguaje objeto} —se llama \ident{tcFn}— que la
teoría aplica a lo que sea. No calcula nada por adelantado: es una expresión del lenguaje, como
$\obj{\sigma}$ o $\obj{+}$.

\begin{center}\small
\begin{tabularx}{\linewidth}{@{}l >{\raggedright\arraybackslash}X >{\raggedright\arraybackslash}X@{}}
\toprule
 & $\codigo{\cdot}$ (\ident{formCode}, \ident{termCode}) & $\dotado{\cdot}$ (\ident{tcFn}) \\
\midrule
quién lo calcula & nosotros, en la metateoría & la teoría, dentro del lenguaje \\
qué admite como argumento & una fórmula o término \textbf{concretos} & cualquier término, incluida \textbf{una variable} \\
qué es el resultado & un término cerrado, escrito & una expresión con variables libres \\
¿se puede poner bajo un $\forall$? & \textbf{no} útilmente & \textbf{sí} \\
\bottomrule
\end{tabularx}
\end{center}

\noindent La cuarta fila es la que importa, y se ve en un ejemplo. Ésta es una de las ecuaciones que
caracterizan a \ident{tcFn}, y dice, para \emph{todo} $x$, cómo se relaciona el código de
$\obj{\sigma}x$ con el de $x$:
\[ \forall x.\quad \dotado{(\obj{\sigma}x)} \;\doteq\; \langle 1,\ \codigo{\obj{\sigma}},\ [\,\dotado{x}\,]\rangle \]
\selee{para todo $x$, el código del sucesor de $x$ es la lista formada por la etiqueta de
aplicación, el código del símbolo $\sigma$, y el código de $x$.}

\begin{muro}[Por qué ahí no se puede escribir $\codigo{x}$]
Se podría \emph{teclear}, y ahí está la trampa: $\codigo{x}$ es legal, pero significa «el código del
término \emph{la variable~$x$}» — un código concreto, cerrado, de un objeto sintáctico. \textbf{No}
significa «el código de lo que valga $x$», que es lo que la fórmula necesita. El punto es la única
notación que puede ir bajo un cuantificador diciendo lo segundo.
\end{muro}

\begin{leccion}[Y una asimetría que se cobra cara]
Como $\dotado{\cdot}$ es un símbolo de función, la teoría demuestra que \textbf{respeta la
igualdad}: de $a \doteq b$ se sigue $\dotado{a} \doteq \dotado{b}$. $\codigo{\cdot}$ no tiene nada
parecido, porque no es un símbolo del lenguaje: es una operación sobre \emph{sintaxis}.

Y ahí está el precio. Un símbolo de función sólo puede depender del \textbf{valor} de su argumento,
nunca de cómo se escribió ese valor. Pero un mismo número es a la vez un numeral y una lista —
$\obj{\mathrm{cons}}(\mathrm{nil},\mathrm{nil})$ y $\obj{\sigma\sigma}\obj{0}$ son el mismo número—,
así que exigirle a $\dotado{\cdot}$ que devuelva un código distinto según «cómo se escribió» es
pedirle lo imposible. Eso hizo la teoría inconsistente durante meses, y es el episodio central de la
Parte~\ref{parte:noescrito}.
\end{leccion}

\section{Cómo se lee la chapa}

De todo lo anterior sale el instrumento que este libro usa en cada fragmento de código y en cada
enunciado: una \textbf{chapa de dos casillas} que dice quién afirma y sobre qué. No es un adorno —
es la distinción cuya violación costó 31 módulos, y \textbf{no es binaria}: son dos ejes
independientes.

%% La leyenda de niveles. Va ANTES del índice: PLAN-LIBRO.md §2.5.
\section*{2 · Los dos ejes: quién afirma, y sobre qué}

Con esas palabras ya se puede plantear la pregunta que vuelve en cada línea:
\emph{¿esto es lenguaje objeto o lenguaje meta?} En un libro de metamatemática se
plantea en cada línea, y responderla mal es caro: la inconsistencia que cuenta la
Parte~IV fue exactamente un error de categoría. Por eso aquí el nivel \textbf{se imprime}.

No es una distinción binaria. Hay dos ejes independientes, y cada fragmento de código y cada
enunciado llevan una chapa con los dos:

\noindent\textbf{Eje 1 — ¿quién afirma?} Son cuatro relaciones de derivabilidad distintas, y
\emph{de tipos distintos}: lo que las separa es si arrastran contexto y si son recursivamente
enumerables.

\begin{center}\footnotesize\setlength{\tabcolsep}{4pt}
\begin{tabularx}{\linewidth}{@{}l >{\ttfamily\scriptsize}l cc >{\raggedright\arraybackslash}X@{}}
\toprule
chapa & tipo en Lean & ctx. & r.e. & papel \\
\midrule
\colorbox{nvLean}{\textcolor{white}{\sffamily\scriptsize\bfseries\strut\,Lean\,}} &
  \normalfont--- & --- & --- & la metateoría: un teorema \emph{sobre} el sistema \\
\colorbox{nvDer}{\textcolor{white}{\sffamily\scriptsize\bfseries\strut\,$\Gamma\vdash$\,}} &
  List Formula → Formula → Prop & sí & \textbf{no} & el cálculo $\omega$; sólido en $\mathbb{N}$ \\
\colorbox{nvPrf0}{\textcolor{white}{\sffamily\scriptsize\bfseries\strut\,Prf$_0$\,}} &
  Formula → Prop & no & sí & Hilbert \textbf{intuicionista} \\
\colorbox{nvPrf}{\textcolor{white}{\sffamily\scriptsize\bfseries\strut\,Prf\,}} &
  Formula → Prop & no & sí & Hilbert clásico: \textbf{el que se aritmetiza} \\
\colorbox{nvPrfH}{\textcolor{white}{\sffamily\scriptsize\bfseries\strut\,$\Gamma\,$PrfH\,}} &
  List Formula → Formula → Prop & sí & sí & variante con contexto; ahí vive la deducción \\
\bottomrule
\end{tabularx}
\end{center}

\noindent\textbf{Eje 2 — ¿sobre qué?} Cinco valores.

\begin{center}\small
\begin{tabularx}{\linewidth}{@{}l >{\raggedright\arraybackslash}X@{}}
\toprule
\colorbox{nvLean!18}{\textcolor{nvLean}{\sffamily\scriptsize\bfseries\strut\,Lean\,}} &
  objetos nativos: \ident{Nat}, listas, recursión estructural \\
\colorbox{nvLean!18}{\textcolor{nvLean}{\sffamily\scriptsize\bfseries\strut\,sintaxis\,}} &
  la representación en Lean de la sintaxis: \ident{Term}, \ident{Formula} \\
\colorbox{nvObj!18}{\textcolor{nvObj}{\sffamily\scriptsize\bfseries\strut\,objeto\,}} &
  enunciados de la teoría aritmética: $\obj{+}$, $\obj{\cdot}$, $\obj{\sigma}$, sus axiomas \\
\colorbox{nvCod!18}{\textcolor{nvCod}{\sffamily\scriptsize\bfseries\strut\,código\,}} &
  términos objeto que \emph{denotan} sintaxis: $\cod{\codigo{\varphi}}$, \ident{numeral} \\
\colorbox{nvProv!18}{\textcolor{nvProv}{\sffamily\scriptsize\bfseries\strut\,Prov\,}} &
  dentro del predicado de demostrabilidad: la teoría hablando de sí misma \\
\bottomrule
\end{tabularx}
\end{center}

\noindent \textbf{Los puentes, y el que no existe.} Las cuatro no son intercambiables, y la
dirección importa:
\[
  \Prf_0\,\varphi \;\longrightarrow\; \Prf\,\varphi \;\longrightarrow\;
  \bigl(\forall \Gamma,\ \mathsf{PrfH}\ \Gamma\ \varphi\bigr),
  \qquad
  \Prf\,\varphi \;\longrightarrow\; \text{\ident{axioms}} \vdash \varphi .
\]
El recíproco $\text{\ident{axioms}} \vdash \varphi \to \Prf\,\varphi$ \textbf{es falso, y tiene que
serlo}: si valiera, $\vdash$ sería r.e.\ y Tarski cerraría el paso. Ésa es la asimetría que obliga a
tener dos cálculos en vez de uno, y por eso la chapa distingue en cuál se afirma.
Además \ident{prf0_to_derives} no usa \emph{ninguna} meta-regla $\omega$ y \ident{prf_to_derives}
usa \ident{dne} exactamente una vez: la diferencia entre los dos puentes es, medida,
$\{\ident{FOL.MetaRules.dne}\}$.

\medskip\noindent Tres ejemplos que ocupan tres casillas distintas y que conviene comparar:
\ident{two_mul_consN} afirma \emph{en Lean} sobre \ident{Nat};
\ident{prf_formCode_numeral} afirma \emph{en Prf} sobre \textbf{códigos};
\ident{hFN} afirma lo mismo, pero \emph{en} $\vdash$.
Y \ident{pcc_dot_cons} afirma en \ident{Prf} \textbf{dentro de} \ident{Prov} — cuatro niveles de
distancia entre el primero y el último.

\begin{leccion}[Por qué la chapa es un control, no un adorno]
La cuarta fila del segundo eje es la que se cobró 31 módulos. \ident{tcFn} —«el código del
código»— es una operación sobre \textbf{sintaxis} que se declaró como función \textbf{objeto}, y
una función objeto sólo puede depender de \textbf{valores}. Con la chapa puesta, la mezcla se ve
antes de escribirla.
\end{leccion}


\begin{leccion}[La distinción que hay que llevarse de aquí]
$\Prov$ es una fórmula del lenguaje objeto; \defterm{provable} quiere decir que la propia teoría
demuestra algo, no que lo comprobemos nosotros desde fuera. La diferencia entre comprobar una cosa y
que la teoría la demuestre \emph{de sí misma} es toda la distancia que hay entre la primera condición
de derivabilidad y la tercera — es decir, entre lo que está hecho y lo que falta.
\end{leccion}

\chapter{Lo mínimo de Lean: cuatro palabras}
\label{cap:lean}

Este libro está lleno de código Lean 4 y no es un libro sobre Lean. Para leerlo hacen falta cuatro
palabras. Con ellas se entiende cada fragmento que aparece; sin ellas, ninguno. Todo lo demás
—tácticas, elaboración, instancias— es maquinaria que el lector puede ignorar sin perder el hilo.

\section{\ident{inductive}: cómo se dice «esto es todo lo que hay»}

Un tipo inductivo se declara enumerando sus \emph{constructores}, y esa enumeración es exhaustiva:
no hay más términos que los que se obtienen aplicándolos.

\leanfrag{fol_Term}

\noindent Un término es una variable o la aplicación de un símbolo de función a una lista de
términos. Nada más. De esa declaración salen gratis tres cosas que en la teoría objeto cuestan
axiomas:

\begin{center}\small
\begin{tabularx}{\linewidth}{@{}l >{\raggedright\arraybackslash}X@{}}
\toprule
constructores disjuntos & una variable nunca es una aplicación \\
constructores inyectivos & si dos aplicaciones son iguales, coinciden símbolo y argumentos \\
inducción estructural & para demostrar algo de todo término basta hacerlo caso por constructor \\
\bottomrule
\end{tabularx}
\end{center}

\begin{leccion}[Lo que regala el compilador, en primer orden se paga]
Esas tres cosas son exactamente lo que la parte~\ref{parte:aritmetizacion} tiene que \emph{demostrar} cuando reconstruye
listas y árboles sintácticos dentro de la aritmética, donde no hay tipos inductivos: hay números.
Ver hecho a mano lo que aquí es una línea enseña qué había realmente dentro de esa comodidad.
\end{leccion}

\section{Las tres palabras que son tres categorías morales}

Las otras tres —\ident{def}, \ident{theorem}, \ident{axiom}— parecen tres formas de escribir lo
mismo. No lo son. Son tres compromisos distintos, y la diferencia entre ellos organiza el proyecto
entero.

\subsection*{\ident{def} — nombrar no cuesta nada}

Una definición introduce una abreviatura. No afirma nada y no puede ser falsa: como mucho, puede ser
inútil.

\leanfrag{consN}

\subsection*{\ident{theorem} — se gana}

Un teorema afirma algo, y sólo entra en el sistema si el núcleo de Lean acepta su demostración.

\leanfrag{two_mul_consN}

\subsection*{\ident{axiom} — es una deuda}

Un axioma afirma algo \textbf{sin demostración}. Lean lo acepta y lo apunta. Es legítimo cuando lo
que se postula es una regla del juego o un principio que la teoría de destino tiene por definición
— y es peligroso siempre, porque un axioma falso hace demostrable cualquier cosa.

\leanfrag{fol_gen}

\begin{muro}[Un axioma falso ya ha pasado aquí, dos veces]
En \modulo{FOL} hubo un \ident{axiom} cuyo enunciado era literalmente falso, con contraejemplo de
tres líneas; hoy es un teorema con el enunciado corregido. Y en la teoría objeto hubo otro que hacía
la teoría inconsistente durante meses. Los dos episodios están contados en este libro, y los dos
empiezan igual: alguien escribió \ident{axiom} donde debía haber escrito una pregunta.
\end{muro}

\section{Por eso hay una línea gris debajo de cada teorema}

Si un axioma es una deuda, lo mínimo es llevar la cuenta. Lean lo permite: dado cualquier teorema,
se puede preguntar de qué axiomas depende — no los que están declarados en el fichero, sino los que
su demostración usa realmente, siguiendo toda la cadena. Esa lista es el \defterm{footprint}, y en
este libro se imprime debajo de cada enunciado citado:

\footprint{\fpi{propext}, \fpi{Classical.choice}, \fpi{Quot.sound}}

\begin{leccion}[Cómo se lee]
Los tres de arriba son la base del propio Lean y salen prácticamente en todo. Lo interesante es lo
que aparece \emph{además}: si un teorema cita un esquema de inducción, es que lo necesita; si no lo
cita, es que no. Cuando el footprint no añade nada a la base ya aceptada se dice que el resultado es
\defterm{net-0}. Y cuando aparece algo que no debería, hay un problema — así se descubrió la
inconsistencia de la parte~\ref{parte:noescrito}.
\end{leccion}

\section{Dos palabras más, que no son de Lean sino de este proyecto}

Aparecen a menudo y conviene fijarlas aquí. Un \defterm{sondeo} es un experimento compilado que
contesta una pregunta —muchas veces respondiendo \emph{que no}— y que vive fuera de la compilación
principal. La \defterm{cuarentena} es el directorio donde se apartaron los módulos que dependían de
la inconsistencia; hoy está vacío, y el episodio es la Parte~\ref{parte:noescrito}.

\begin{leccion}[Todo lo anterior en una línea]
\ident{def} nombra, \ident{theorem} se gana, \ident{axiom} se debe. El resto de este libro consiste,
básicamente, en pagar deudas y en anotar las que no se han podido pagar todavía.
\end{leccion}


% =============================================================================
\part{El terreno}
\label{parte:terreno}

\noindent \emph{¿Con qué se escribe, y qué se afirma?} La primera y la segunda de las cinco piezas: el lenguaje —con sus dos cálculos, por la razón que dio el capítulo~\ref{cap:lenguajes}— y la teoría aritmética sobre la que se demuestra todo lo demás.
% =============================================================================
\include{capitulos/cap-kernel-fol}             % 5 ✅
\chapter{¿Cuánta metateoría hace falta?}
\label{cap:objeto-y-meta}

El capítulo~\ref{cap:kernel} dejó los dos cálculos y el puente entre ellos; el
capítulo~\ref{cap:qpp} dirá qué se afirma con ellos. Entre los dos queda una pregunta que casi
ningún libro de incompletitud hace explícita, y que aquí no se puede eludir porque la metateoría
es una máquina concreta: \emph{¿cuánta fuerza lógica estamos gastando para hablar de la teoría
objeto, y cuánta hace falta de verdad?}

La pregunta no es ociosa. Gödel escribió su metateoría en prosa, pero deliberadamente
\textbf{finitaria}, para que Hilbert la aceptara: no valía razonar sobre las pruebas con más
recursos que los que se estaban poniendo en duda. Aquí la metateoría es Lean~4, que es
enormemente más fuerte que eso. Si nadie mide cuánto se usa, el libro estaría demostrando la
incompletitud de una aritmética débil desde una torre de la que no sabe nada. Este capítulo mide.

\begin{leccion}[La asimetría que hay que decir en voz alta]
La metateoría \textbf{no tiene que ser más fuerte} que la teoría de la que habla. Sólo tiene que
poder hablar de \emph{sintaxis finita}: fórmulas, listas de fórmulas, y la comprobación mecánica de
que una lista es una derivación. Es una exigencia de \emph{expresividad}, no de fuerza. Esto es lo
que hace posible el proyecto entero — y lo que hace que la pregunta «¿cuánta metateoría?» tenga
una respuesta pequeña.
\end{leccion}

\section{Lean no es la aritmética de Peano, y por tres razones}

Conviene descartar primero la respuesta cómoda. Sería elegante que la metateoría de este proyecto
fuera exactamente PA —la aritmética de primer orden con inducción—: el libro se cerraría sobre sí
mismo. No lo es, y por bastante margen.

\textbf{Primero, la inducción de Lean no es un esquema.} La de PA sí: un axioma \emph{por cada}
fórmula del lenguaje aritmético, infinitas numerables, todas cuantificando sólo sobre números. Eso
es un \defterm{esquema de axiomas}\nivelterm{objeto}. El principio de inducción que Lean genera
para un tipo inductivo es \textbf{un solo enunciado} que cuantifica sobre un motivo
\ident{C : Formula → Sort u} — es decir, sobre \emph{todos} los predicados de fórmulas. Eso es
inducción de segundo orden, y no hay manera de escribirla como esquema numerable.

Y no es una diferencia de manual: el proyecto la usa así de verdad. La demostración de que el
código de una fórmula es igual a su numeral va por recursión meta sobre \ident{Formula} con el
motivo $\varphi \mapsto \Prf(\text{\ident{formCode}}\ \varphi \doteq \text{\ident{numeral}}
(\text{\ident{codeNat}}\ \varphi))$, y $\Prf$ es a su vez un \ident{Prop} inductivo — un motivo que
en PA no se puede ni enunciar sin haber aritmetizado antes $\Prf$.

\leanfrag{prf_formCode_numeral}

\textbf{Segundo, los universos.} Lean tiene una jerarquía infinita de universos con \ident{Prop}
impredicativo. Su teoría de tipos se interpreta en \textsc{zfc} más una cantidad numerable de
cardinales inaccesibles, y dentro de Lean puede construirse un modelo de \textsc{zfc}; luego Lean
demuestra $\mathrm{Con}(\textsc{zfc})$, y con ella $\mathrm{Con}(\mathrm{PA})$, que PA no
demuestra.\footnote{M.~Carneiro, \emph{The Type Theory of Lean}, tesis de máster, Carnegie Mellon
University, 2019. Es la referencia estándar para el modelo; el enunciado exacto conviene
consultarlo allí.}

\textbf{Tercero, los axiomas.} \ident{propext}, \ident{Classical.choice} y \ident{Quot.sound} están
en el \defterm{footprint}\nivelterm{proyecto} de prácticamente todo lo que este libro cita: son la
base sancionada del proyecto. Ninguno de los tres tiene contrapartida aritmética.

\section{Cuánto de esa fuerza se usa: dos niveles, y ninguno más}

Que Lean sea fuerte no dice nada sobre cuánto se gasta. Eso se mide, y el resultado es
sorprendentemente pequeño.

\begin{center}\small
\begin{tabular}{lll}
\toprule
qué & tipo & nivel \\
\midrule
\ident{Term}, \ident{Formula}, \ident{Pos} & \ident{Type} & \ident{Sort 1} \\
\ident{Derives}, \ident{Prf}, \ident{PrfH} & \ident{Prop} & \ident{Sort 0} \\
\bottomrule
\end{tabular}
\end{center}

En todo el árbol de \modulo{ROBINSON_PlusPlus/} y en los módulos de \modulo{FOL} que la cadena de
\ident{import} alcanza no hay \textbf{ni un} \ident{Sort u}, \textbf{ni un} \ident{Type 1},
\textbf{ni una} declaración \ident{universe}. El único \ident{Type u} del árbol vive en
\modulo{FOL/Completeness.lean}, que —como dijo el capítulo~\ref{cap:kernel}— ningún módulo importa.

Aunque el recursor de \ident{Formula} sea polimórfico en universos, el proyecto lo instancia en
exactamente dos sitios: motivo en \ident{Prop} cuando \emph{demuestra} por inducción, y motivo en
\ident{Type} cuando \emph{define} por recursión (\ident{codeNat}, \ident{formCode},
\ident{liftTerm}, \ident{substFormula}). Ninguna inducción concreta necesita el polimorfismo: cada
una vive en un nivel fijo.

\begin{observacion}
Y algo más fuerte: \textbf{no se cuantifica sobre \ident{Prop} en ninguna parte} —ni
\ident{(P : Prop)} ni \ident{{P : Prop}}—, luego la impredicatividad de \ident{Prop}, que es de
donde sale buena parte de la fuerza de Lean, \textbf{no se usa}. Dato relacionado: hay un solo
\ident{termination_by} en todo el árbol (\modulo{ROBINSON_PlusPlus/Meta/CodeDecode.lean}); todo lo
demás es recursión estructural.
\end{observacion}

\begin{muro}[Qué mide exactamente esto]
Todas las cifras de esta sección están medidas sobre el \textbf{texto} del árbol, con búsquedas
literales. Un \printaxioms\ sobre el entorno compilado puede destapar usos implícitos que una
búsqueda no ve —una táctica que instancia algo en un universo alto, por ejemplo—. El resultado hay
que leerlo como «no hay ningún uso \emph{escrito}», que es menos de lo que uno querría y bastante
más de lo que suele decirse.
\end{muro}

\section{La traducción a una metateoría aritmética}

Dos niveles, y cada uno tiene contrapartida conocida en aritmética:

\begin{itemize}
\item motivo en \ident{Type} —definir una función recursiva sobre sintaxis— se traduce por
  \defterm{recursión de curso de valores}\nivelterm{meta}: la función $\beta$ de Gödel, que permite
  codificar una sucesión finita en un número y así definir por recursión dentro de la aritmética.
  En este proyecto la capa \ident{lenc}/\ident{nthc} ya se construye \emph{para otra cosa}, y es
  exactamente ese instrumental;
\item motivo en \ident{Prop} —probar por inducción— se traduce por el esquema de inducción de la
  metateoría.
\end{itemize}

La fuerza necesaria es del orden de $\mathrm{I}\Sigma_1$: inducción restringida a fórmulas
$\Sigma_1$. Y aquí se cierra un círculo que merece decirse: \doc{AXIOMS.md}~§1.1 argumenta que la
\emph{teoría objeto} necesita inducción de ese mismo perfil para D2 y D3. Teoría y metateoría
piden lo mismo. No es casualidad: las dos tienen que hacer el mismo trabajo, hablar de objetos
sintácticos finitos.

\section{\texorpdfstring{\texttt{Classical}}{Classical}: dos usos, y no son de la misma clase}

En todo \modulo{ROBINSON_PlusPlus/} hay \textbf{dos} usos explícitos de lógica clásica. Dos. Y la
diferencia entre ellos es el resultado más interesante de este capítulo.

\textbf{(a) Eliminable.} En \modulo{ROBINSON_PlusPlus/Full/PrimeFactor.lean} se saca
$\exists a,\ a \mid n \wedge a \neq 1 \wedge a \neq n$ por reducción al absurdo desde
«$n$ no es primo».

\leanfrag{exists_prime_factor}

Pero todos los divisores están acotados por $n$, y la divisibilidad, la igualdad y «ser primo» son
decidibles sobre \ident{Nat}. Es una \textbf{búsqueda acotada}: constructiva de manual.
\ident{Decidable.byContradiction} la sustituye sin perder nada.

\textbf{(b) No eliminable — y está identificado.} El otro está en la reducción que descarga la
hipótesis de reflexión:

\leanfrag{Reflects}
\leanfrag{reflects_of_omega}

Tras el \ident{intro}, el objetivo es $\Prf\varphi$ con la hipótesis $\neg\,\Prf\varphi$: o sea,
$\neg\neg\,\Prf\varphi \to \Prf\varphi$. Y $\Prf\varphi$ es $\Sigma_1$ —existe una derivación
finita, y comprobar que una lista dada es una derivación es decidible: ése es justo el trabajo del
verificador del capítulo~\ref{cap:verificador}—.

\begin{leccion}[Eso tiene nombre, y es el más débil que sirve]
Eliminar la doble negación sobre un enunciado $\Sigma_1$ de matriz decidible es el
\defterm{principio de Markov}\nivelterm{meta}. No es «lógica clásica» a secas: es el principio
\textbf{más débil} que separa la matemática constructiva intuicionista de la recursiva rusa. Poder
decir «un solo uso esencial, y es Markov» no es una nota de limpieza — es un resultado sobre el
proyecto.
\end{leccion}

El \ident{Classical.choice} que aparece en los footprints es otra cosa: entra casi todo por
instancias \ident{Decidable} y por tácticas. En
\modulo{ROBINSON_PlusPlus/Meta/ChainDecode.lean} se construyen a mano \ident{DecidableEq} para
\ident{Term} y \ident{Formula} —el \ident{deriving} automático no cubre el anidamiento
\ident{func : String → List Term}—, así que hay decidibilidad real disponible y ese uso es en
principio prescindible.

\begin{muro}[Lo que un \printaxioms\ no distingue]
Los footprints son la moneda del proyecto, pero \ident{Classical.choice} en un \printaxioms\ sale
igual venga del Markov esencial o de una táctica perezosa. Y esa distinción es justo la que hace
falta. La sección siguiente cuenta que la herramienta que la hace ya existe.
\end{muro}

\section{El experimento ya está hecho, y en el mismo ecosistema}

Sustituir la maquinaria meta de Lean por una metateoría aritmética es un objetivo declarado del
autor: usar su proyecto \modulo{Peano} —con alguna extensión— en lugar de Lean. La pregunta
natural es si eso es una aspiración o un plan. Resulta que es un plan, porque la parte difícil ya
se hizo allí.

\modulo{Peano} lleva desde 2026 una \textbf{guarda de compilación}: un comando que, en cada
\ident{lake} \ident{build}, comprueba que los teoremas vigilados no dependen de
\ident{Classical.choice}. Si alguien introduce lógica clásica por accidente, la compilación falla.

\leanfrag{peano_assert_constructive}

Hoy hay \textbf{1425} invocaciones activas de ese comando. Y —el detalle que lo hace confiable— el
fichero conserva, comentada, una invocación sobre un teorema que \emph{sí} usa elección, para poder
verificar que la guarda detecta lo que dice detectar. Es un control positivo, exactamente la
disciplina que este libro se impone a sí mismo en \doc{PLAN-LIBRO.md}~§2.

\begin{leccion}[Lo que eso demuestra para el plan de sustitución]
La técnica es \textbf{portable tal cual}: es un fichero, no una metodología por inventar.
Aplicarla aquí daría la separación que la sección anterior pide —Markov esencial frente a
\ident{Decidable} incidental— como \emph{control de compilación} y no como auditoría manual. Es la
doctrina de este libro aplicada a sí mismo: convertir un principio en una máquina.
\end{leccion}

Y hay un dato que vale más que el método. La función $\beta$ de Gödel —la pieza que la sección
anterior identifica como necesaria para que una metateoría aritmética hospede la recursión de curso
de valores— figuraba en la lista inicial de deuda clásica de \modulo{Peano}. Ya no: sus doce
declaraciones principales están hoy bajo la guarda, y la reparación consistió en reemplazar la
elección por \textbf{búsqueda acotada} — la misma técnica que arriba se propone para el uso
eliminable de este proyecto. La única excepción que \modulo{Peano} mantiene, documentada y
aceptada, es un módulo cuyo cometido es cuantificar sobre sistemas de Peano \emph{abstractos}
cualesquiera: metateoría genuinamente no constructiva, sin nada que enumerar ni acotar.

\begin{observacion}
Nótese qué acaba de pasar aquí, porque es un caso de \doc{DOCSTRINGS-NO-FIABLES.md} en vivo. La
lista de deuda clásica que la documentación de \modulo{Peano} exhibe sigue nombrando el módulo de
la función $\beta$; el árbol dice que se resolvió. La documentación es el inventario de partida, no
el estado. Lo que decide es la medición.
\end{observacion}

\section{Qué transferiría y qué no}

Suponiendo hecha la sustitución, conviene decir por adelantado cuál sería su alcance, porque no es
total y la mitad que falta no es un defecto.

\begin{center}\small
\begin{tabularx}{\linewidth}{@{}>{\raggedright\arraybackslash}p{0.22\linewidth}
                              >{\raggedright\arraybackslash}X@{}}
\toprule
qué & transfiere a una metateoría de fuerza $\mathrm{I}\Sigma_1$ \\
\midrule
todo $\Prf$ & \textbf{sí, entero}: es finitario, r.e.\ y por recursión estructural sobre objetos
finitos. D1, D2, el punto fijo y Gödel~II sobre Q\texttt{++} caben ahí \\
$\derives$ con \ident{gen} & \textbf{no}: la \defnot{$\omega$-regla} tiene premisas infinitas, y
ninguna aritmética de primer orden puede hospedarla \\
\bottomrule
\end{tabularx}
\end{center}

O sea: una metateoría-Peano reproduciría la mitad $\Prf$ del proyecto y no la mitad $\derives$. No
es una limitación nueva — es exactamente la misma razón por la que el capítulo~\ref{cap:kernel}
tenía que presentar \emph{dos} cálculos y no uno. Pero la sustitución habría que enunciarla como
parcial, y decir cuál mitad, que es lo que casi nunca se dice.
          % 6 ✅
\include{capitulos/cap-robinson-qpp}           % 7 ✅

% =============================================================================
\part{Aritmetización}
\label{parte:aritmetizacion}

\noindent \emph{¿Cómo consigue una teoría de números hablar de fórmulas?} Las tres piezas que faltan: la codificación, el verificador y la primera condición de derivabilidad. Al final de esta parte, $\Prov$ existe y está demostrado que dice lo que pretende.
% =============================================================================
\chapter{Codificar estructuras como números}
\label{cap:estructuras}

En una teoría de primer orden no hay tipos inductivos. No hay listas, no hay árboles, no hay
fórmulas: hay números, funciones y predicados. Y sin embargo Gödel necesita que la teoría hable de
listas y de árboles sintácticos. Este capítulo es el puente.

\section{Un par es un número}

La pieza de partida es el emparejamiento de Cantor, que el capítulo~\ref{cap:qpp} ya demostró
biyectivo dentro de la teoría. Con él, un par ordenado deja de ser una estructura y pasa a ser un
número.

\begin{matematica}[De pares a listas]
\[
  \langle x,y\rangle \;=\; \frac{(x{+}y)(x{+}y{+}1)}{2} + y,
  \qquad
  \obj{\mathrm{nil}} \;=\; \obj{0},
  \qquad
  \obj{\mathrm{cons}}(h,t) \;=\; \langle h,\ \obj{\sigma}t\rangle .
\]
El \(\obj{\sigma}\) de la cola no es adorno: hace que ningún \(\obj{\mathrm{cons}}\) valga
\(\obj{0}\), y por tanto que \(\obj{\mathrm{nil}}\) no sea confundible con una lista no vacía.
\end{matematica}

La pertenencia se caracteriza con un solo axioma, y es el que hace que una lista se pueda recorrer:

\leanfrag{ax_in_cons}

\begin{leccion}[Lo que se gana y lo que se paga]
Se gana que \emph{toda} la teoría de listas viva dentro de la aritmética, sin ampliar la lógica. Se
paga que \textbf{un número es varias cosas a la vez}: el mismo valor es un par, un sucesor y una
lista, y nada en el número dice cuál se pretendía. Esa ambigüedad es exactamente el material del
capítulo~\ref{cap:numerales}.
\end{leccion}

\section{Los accesos: símbolos opacos con axiomas}

Para recorrer listas hacen falta cabeza, cola, longitud e índice. Y aquí aparece una decisión de
diseño que conviene ver, porque explica buena parte del recuento de axiomas: estos símbolos
\textbf{no se definen}, se declaran opacos y se fijan con axiomas.

\leanfrag{carc}
\leanfrag{lenc}
\leanfrag{nthc}

\noindent Un \ident{Term.func "carc" [l]} es literalmente eso: un símbolo de función del lenguaje
objeto, sin contenido. Lo que le da significado son sus ecuaciones —\(\obj{\mathrm{carc}}(h::t) =
h\), \(\obj{\mathrm{lenc}}(\mathrm{nil}) = \obj{0}\), y las demás—, que viven en la lista de
axiomas.

\begin{leccion}[Símbolo opaco frente a definición]
Un símbolo \emph{definido} —como \(\obj{\mathrm{cons}}\), que es
\(\obj{\mathrm{div2}}\) del polinomio de Cantor— no cuesta axiomas: su comportamiento se deduce.
Un símbolo \emph{opaco} cuesta tantos axiomas como ecuaciones haga falta fijar. Es el mismo criterio
que retiró \ident{ax22_cantor_proj_exists} y \ident{ax23_cantor_proj_uniq} en 2026: convertir dos
símbolos opacos en definiciones eliminó sus axiomas de golpe.
\end{leccion}

\section{La inducción estructural, recuperada}

Queda una objeción de fondo. En Lean, \ident{List} viene con su principio de inducción; en la teoría
objeto no hay tal cosa. ¿Cómo se demuestra algo «para toda lista»?

\begin{matematica}[De la inducción numérica a la estructural]
Se define la longitud \(\obj{\mathrm{lenc}}\) como función aritmética y se induce \emph{sobre ella}:
si la propiedad vale para la lista vacía, y se conserva al añadir un elemento, entonces vale para
toda lista — porque toda lista tiene una longitud, y sobre los números sí hay inducción.
\end{matematica}

Esto exige inducción, luego \textbf{no está disponible en \modulo{Minimal/}}: ahí los dos hechos
inductivos sobre listas que se necesitan (\ident{ax_C3_concat_assoc} y \ident{ax_L3_in_concat}) se
postulan, y pasan a teoremas en \modulo{Full/}. Es la misma frontera del capítulo anterior, vista
desde las estructuras de datos.

\begin{leccion}[El engarce con los tipos inductivos de Lean]
Lo que este capítulo hace a mano —inyectividad y disyunción de los constructores, inducción
estructural desde la numérica, funciones recursivas por recorrido— es exactamente lo que el
compilador de Lean regala al declarar un \ident{inductive}. Verlo hecho a mano en primer orden
enseña qué hay realmente dentro de esa comodidad.
\end{leccion}
  % 8  ✅
\chapter{Códigos de términos y fórmulas}
\label{cap:codigos}

Con listas dentro de la aritmética, ya se puede codificar sintaxis. Este capítulo construye
\(\codigo{\varphi}\): el término del lenguaje objeto que denota la fórmula \(\varphi\).

\section{Dos codificaciones, y por qué hay dos}

El proyecto tiene una codificación \emph{de juguete} y una \emph{real}, y conviene no confundirlas.

La primera —el nivel B— codifica cadenas de símbolos de un alfabeto finito de doce elementos:

\leanfrag{Sym}
\leanfrag{gNat}
\leanfrag{encode}

\noindent Es la construcción clásica de Gödel, útil para explicar la idea y para el teorema de
inyectividad, pero no sirve para el argumento real: codifica \emph{cadenas}, no el árbol sintáctico
de las fórmulas del kernel.

\section{La codificación real: por constructor}

La segunda codificación —el nivel C— recorre la estructura de \ident{Term} y \ident{Formula} y
asigna una etiqueta a cada constructor.

\leanfrag{termCode}
\leanfrag{formCode}

\begin{matematica}[Las etiquetas]
\[
\begin{array}{llll}
\bot \mapsto 2 & \text{átomo} \mapsto 3 & {\doteq} \mapsto 4 & {\Rightarrow} \mapsto 5\\
\forall \mapsto 6 & \wedge \mapsto 7 & \vee \mapsto 8 & \exists \mapsto 9
\end{array}
\]
y, en los términos, \(0\) para variables y \(1\) para aplicaciones. Un código es una lista cuya
cabeza es la etiqueta y cuya cola son los códigos de las partes.
\end{matematica}

\begin{leccion}[Dónde vive cada cosa]
\ident{formCode} es una función \textbf{de Lean}: toma un valor del tipo \ident{Formula} y devuelve
un valor del tipo \ident{Term}. Su \emph{resultado} es un término del lenguaje objeto. Por eso la
chapa dice \emph{Lean} sobre \emph{código}, y por eso el capítulo de nomenclatura insiste: la
función que calcula el código es meta; el código es objeto.
\end{leccion}

\section{Inyectividad, sin suponer consistencia}

Que la codificación sea inyectiva es lo que la hace utilizable: dos fórmulas distintas tienen
códigos distintos, luego un código determina su fórmula.

\begin{teorema}[La codificación es inyectiva]
Si \(\codigo{\varphi} = \codigo{\psi}\) entonces \(\varphi = \psi\).
\end{teorema}

\leanfrag{formCode_injective}

\begin{leccion}[Un detalle que importa más de lo que parece]
Este teorema se demuestra \textbf{en la metateoría}, por inducción estructural, y \textbf{no supone
la consistencia} de la teoría objeto. Es un hecho sobre los constructores de un tipo inductivo de
Lean, no sobre lo que la aritmética demuestre. Si dependiera de la consistencia, todo el edificio
sería circular.
\end{leccion}

\section{De códigos a demostrabilidad}

Con la codificación se pueden definir, ya en la metateoría, los dos predicados que dan sentido a
todo lo demás:

\leanfrag{IsFormula}
\leanfrag{Provable}

\noindent Y el primer resultado que los conecta:

\leanfrag{provable_formCode_iff}

\begin{muro}[Esto todavía no es Gödel]
\ident{Provable} es un predicado \emph{de Lean} sobre términos: dice que existe una fórmula con ese
código y que la teoría la demuestra. Es correcto y es inútil para la incompletitud, porque la
teoría objeto no puede hablar de él — vive en la metateoría. Lo que hace falta es una \emph{fórmula
del lenguaje objeto} que exprese «existe una demostración». Ése es el capítulo siguiente.
\end{muro}
    % 9  ✅
\include{capitulos/cap-el-verificador}         % 10 ✅
\include{capitulos/cap-representabilidad-d1}   % 11 ✅

% =============================================================================
\part{Los teoremas}
\label{parte:teoremas}

\noindent \emph{Y ahora sí: los dos teoremas de Gödel.} El punto fijo, la mitad demostrada del primero y la que falta, las tres condiciones de derivabilidad y el segundo teorema — módulo la única que sigue sin cerrarse.
% =============================================================================
% \chapter{Autorreferencia}
\label{cap:autorreferencia}

Al final de la parte anterior el edificio está montado: hay una teoría de números, hay una
codificación de la sintaxis, hay un verificador de demostraciones y hay una fórmula $\Prov$ que
dice, dentro del lenguaje, «existe una demostración de la fórmula cuyo código es $x$». Falta el
disparo. Este capítulo contesta la pregunta que lo hace posible: \emph{¿cómo se consigue que una
fórmula concreta hable de sí misma?}

La respuesta no es un truco de lenguaje. Es una construcción, y aquí está entera y compilada.

\section{El mentiroso, y el cambio que lo salva}

La paradoja del mentiroso —«esta frase es falsa»— es vieja y es un callejón. Si la frase es
verdadera, es falsa; si es falsa, es verdadera. Como argumento matemático no sirve de nada: lo
único que demuestra es que «verdadero» no se deja definir así.

Y esa es, literalmente, la mitad del descubrimiento. El teorema de Tarski dice que \textbf{la verdad
aritmética no es definible en el lenguaje de la aritmética}: no hay ninguna fórmula
$\mathrm{Verdad}(x)$ que valga exactamente para los códigos de las sentencias verdaderas. Si Gödel
hubiera necesitado esa fórmula, no habría teorema.

\begin{leccion}[El cambio de una palabra]
Gödel no dice «esta frase es falsa». Dice \textbf{«esta frase no es demostrable»}. Y la diferencia
lo es todo, porque la demostrabilidad \emph{sí} es expresable: una demostración es un objeto
finito, verificarla es una comprobación mecánica, y el capítulo~\ref{cap:verificador} construyó esa
comprobación dentro de la teoría. \emph{Ser demostrable} se puede escribir; \emph{ser verdadero},
no.

Cambiar «falsa» por «no demostrable» convierte una paradoja en un teorema. Todo lo que este libro
lleva construido —los códigos, el verificador, $\Prov$— existe para pagar ese cambio.
\end{leccion}

Queda el otro problema, y es el de este capítulo: ¿cómo se escribe el «esta»?

\section{Autoaplicación: sustituirse el propio código}

En la metateoría la maniobra es sencilla de decir. Sea $\psi$ una fórmula con una variable libre.
Su \emph{autoaplicación} es lo que resulta de sustituir en esa variable el código de la propia
$\psi$:

\leanfrag{selfApp}

\selee{«\ident{selfApp} $\psi$ es $\psi$ con su propio código de Gödel puesto en su variable
libre~0»}

Es una función de la \textbf{metateoría}: toma una fórmula y devuelve otra. No hay nada
autorreferente todavía — $\psi$ no se menciona a sí misma, somos nosotros los que le enchufamos su
código desde fuera. Y así, sola, no sirve: para que la teoría razone sobre la sentencia resultante,
la teoría tiene que poder \emph{hacer} esa sustitución, no verla hecha.

\section{Diagonalizar dentro del lenguaje}

Aquí es donde se cobra la distinción del capítulo~\ref{cap:lenguajes} entre las dos maneras de
decir «el código de». Hay que escribir un término del lenguaje objeto que, aplicado al código de
una fórmula, dé el código de su autoaplicación. Y ese término tiene que operar sobre una
\textbf{variable}, no sobre una fórmula concreta:

\leanfrag{diagTerm}

\selee{«\ident{diagTerm} es: sustituir, en la variable~0 de la fórmula de código \texttt{\#0}, el
código del código de \texttt{\#0}»}

Léase con cuidado, porque en esa línea está todo el capítulo. El tercer argumento de
\ident{substfc} es \texttt{\#0}: la fórmula en la que se sustituye es \emph{la que sea}, dada por su
código. El segundo es $\dotado{\texttt{\#0}}$: lo que se mete dentro es el \textbf{código del
código}, calculado por la teoría. Por eso hace falta $\dotado{\cdot}$ y no valdría $\codigo{\cdot}$
— $\codigo{\texttt{\#0}}$ sería «el código del término \emph{la variable cero}», un número fijo, y
no lo que aquí se necesita.

\begin{leccion}[Por qué se llama diagonalización]
La maniobra es la de Cantor. Se tiene una familia de fórmulas indexada por números —la fórmula de
código $n$— y se construye algo que consulta la casilla $(n,n)$: la fórmula de código $n$
\emph{aplicada a} $n$. La diagonal de la tabla. Cantor la usó para escapar de una enumeración;
Gödel, para no salirse de ella.
\end{leccion}

\section{Que la teoría sepa hacerlo}

Escribir el término no basta: hay que \textbf{demostrar} que hace lo que se dice, y demostrarlo
\emph{dentro de la teoría}. Ese es el enunciado que convierte la idea en matemáticas:

\leanfrag{diag_arith_num}

\selee{«la teoría demuestra que sustituir el código de $\psi$ en \ident{diagTerm} da exactamente el
código de la autoaplicación de $\psi$»}

Nótese quién afirma: la chapa dice $\derives$ sobre \emph{código}. No somos nosotros comprobándolo
caso a caso —eso sería una afirmación de la metateoría—; es la teoría objeto la que lo demuestra,
para el $\psi$ que se le dé. Sin este teorema, $\Prov$ y \ident{diagTerm} serían dos piezas que no
encajan.

\begin{observacion}
La demostración es una cadena de cuatro igualdades encadenadas por transitividad, y las cuatro
están en el fichero. Merece la pena mirarla porque enseña de qué está hecho realmente este
proyecto: dos congruencias (\emph{si los argumentos son iguales, los resultados lo son}), la
ecuación que caracteriza la sustitución sobre códigos, y el puente entre las dos maneras de
escribir un código que la Parte~\ref{parte:noescrito} explica. Nada más profundo que eso — y aun
así costó reconstruirla entera.
\end{observacion}

\section{La sentencia}

Con las piezas anteriores, $G$ se escribe en tres líneas. Primero el predicado: \emph{no ser
demostrable}.

\leanfrag{godelPred}

Después se diagonaliza ese predicado — es decir, se le antepone \ident{diagTerm}, de modo que
$\beta$ dice, de la fórmula de código \texttt{\#0}, que \emph{su autoaplicación} no es demostrable:

\leanfrag{godelBeta}

Y por fin la sentencia: $\beta$ aplicada a su propio código.

\leanfrag{godelCN}

\selee{«$G$ es $\beta$ con el código de $\beta$ sustituido en su variable libre»}

\begin{matematica}[La sentencia de Gödel, en una línea]
\[ G \;=\; \beta(\codigo{\beta}), \qquad\text{donde}\qquad
   \beta(x) \;=\; \lnot\,\Prov\bigl(\mathrm{diag}(x)\bigr). \]
$G$ no contiene ningún «esta»: es una fórmula cerrada, con un código de Gödel concreto, que se
podría imprimir. Lo autorreferente no está en la sentencia — está en que su código aparece dentro
de ella.
\end{matematica}

\section{El punto fijo}

Y ahora el teorema del capítulo. Un \defterm{punto fijo}\nivelterm{objeto} es una sentencia que la
teoría prueba equivalente a lo que se afirma \emph{de su propio código}. Aquí:

\leanfrag{godelCN_fixedpoint}

\selee{«la teoría demuestra que $G$ es equivalente a que $G$ no sea demostrable»}

\begin{leccion}[Lo que hay que mirar de este enunciado]
Dos cosas, y las dos importan más que la fórmula.

\textbf{No tiene hipótesis.} Ni consistencia, ni $\omega$-consistencia, ni reflexión, ni ningún
postulado gödeliano. Es un teorema de la teoría objeto, a secas. Todo lo que viene después
—Gödel~I, Gödel~II— añadirá hipótesis; el punto fijo no necesita ninguna.

\textbf{Y no dice nada sobre la verdad de $G$.} Sólo dice que la teoría demuestra una
\emph{equivalencia}. De ahí todavía no se sigue que $G$ sea indemostrable: para eso hace falta el
argumento del capítulo siguiente, y hace falta suponer que la teoría es consistente.
\end{leccion}

\section{Dos maneras de escribir el mismo código, y un paso de Leibniz}

Queda un detalle que en otro libro sería una nota al pie y aquí es una cicatriz visible. El punto
fijo se demuestra primero en una forma \emph{numeral} —con el código escrito como el numeral de un
número— y sólo después se traslada a la forma que el argumento de Gödel espera:

\leanfrag{godelCN_fixedpoint_N}

El traslado es un único paso de congruencia, y es este:

\leanfrag{provCode_transfer}

\selee{«las dos maneras de escribir $\Prov$ del código de $\varphi$ son equivalentes en la teoría»}

\begin{leccion}[Por qué hay dos, y por qué eso salvó meses de trabajo]
Porque la primera versión de todo esto era \textbf{inconsistente}. El código se escribía como un
árbol, y el axioma que relacionaba ese árbol con su propio código hacía que la teoría demostrara
$\bot$. La reparación —el capítulo~\ref{cap:numerales}— consistió en escribir los códigos como
numerales.

Lo que hizo barata la reparación fue exactamente este lema. \ident{diagTerm}, $\beta$ y la
composición de sustituciones \textbf{no mencionan cómo se escribe un código}, así que se
reutilizaron sin tocarlos; y D1 y el argumento modular de Gödel siguieron valiendo porque
\ident{provCode_transfer} los conecta con un paso de Leibniz. Un cambio en los cimientos que no
obligó a rehacer la planta de arriba: eso no pasa por suerte, pasa porque las piezas estaban
separadas.
\end{leccion}

\section{Lo que \texorpdfstring{$G$}{G} no dice}

Conviene cerrar el capítulo desactivando tres lecturas que la divulgación repite y que este
libro, que tiene la fórmula delante, puede desmentir con precisión.

\begin{muro}[Tres cosas que $G$ no es]
\textbf{$G$ no dice «yo soy indemostrable».} Dice, literalmente, que un cierto número no está en
cierta relación aritmética con otros números. Que ese número sea el código de la propia $G$ es un
teorema —el punto fijo—, no una lectura.

\textbf{$G$ no es indemostrable «en general».} Es indemostrable \emph{en esta teoría}, y la teoría
está fijada por la lista \ident{axioms}. Cambiar la lista cambia $\Prov$, y cambiar $\Prov$ cambia
$G$: es otra sentencia y otro teorema. El capítulo~\ref{cap:lenguajes} lo dijo con los cuatro
estratos; aquí se ve por qué.

\textbf{$G$ no es una curiosidad autorreferente.} La construcción es uniforme: el mismo argumento
produce un punto fijo para \emph{cualquier} predicado con una variable libre. Que se aplique a
«no demostrable» es una elección; el mecanismo no sabe de qué habla.
\end{muro}

Con el punto fijo en la mano, el primer teorema de Gödel es corto. Es el capítulo siguiente — y la
sorpresa es que sólo la mitad sale gratis.
          % 12
% \chapter{El primer teorema, y la mitad que sale gratis}
\label{cap:godel-i}

Con el punto fijo demostrado, el primer teorema de Gödel ocupa seis líneas de Lean y no necesita
ninguna idea nueva. Este capítulo las lee despacio, porque en ellas se ve exactamente
\emph{qué} se usa — y, sobre todo, qué no.

Y luego dice lo que casi ningún libro dice en el mismo sitio: que de las dos mitades del teorema,
en esta formalización sólo una está cerrada.

\section{El argumento}

Supóngase que $G$ \emph{sí} fuera demostrable. Entonces:

\begin{matematica}[Cuatro pasos y una contradicción]
\begin{enumerate}
\item Si $G$ tiene una demostración finitaria, la teoría lo \emph{sabe}: por D1,
  $\text{\ident{axioms}} \derives \Prov(\codigo{G})$. Eso es exactamente lo que el
  capítulo~\ref{cap:d1} demostró.
\item De $\Prf G$ se pasa a $\text{\ident{axioms}} \derives G$ por el puente entre los dos cálculos
  (capítulo~\ref{cap:kernel}).
\item El punto fijo, leído de izquierda a derecha, dice
  $G \Rightarrow \lnot\Prov(\codigo{G})$. Con (2), la teoría demuestra
  $\lnot\Prov(\codigo{G})$.
\item Junto a (1): la teoría demuestra una fórmula y su negación. Luego demuestra $\bot$.
\end{enumerate}
Si la teoría es consistente, (4) es imposible. Luego $G$ no es demostrable. $\square$
\end{matematica}

Obsérvese que no se ha hablado ni una vez de si $G$ es \emph{verdadera}. El argumento entero es
sintáctico: entra una demostración hipotética y sale una contradicción.

\section{El mismo argumento, en Lean}

Aquí está, entero, sin recortes. Es uno de los pocos sitios de este libro donde merece la pena
imprimir la prueba completa y no sólo el enunciado, porque las cinco líneas del cuerpo son las
cuatro del argumento de arriba más la contradicción:

\leanfrag{goedel_first_unprovable}

\selee{«si la teoría es consistente y $G$ es un punto fijo de $\lnot\Prov$, entonces $G$ no tiene
demostración finitaria»}

Léase la correspondencia línea a línea: \ident{repr_pos'} es D1; \ident{prf_to_derives} es el
puente entre los cálculos; \ident{and_elim_left} extrae la dirección del bicondicional que hace
falta; y la hipótesis \texttt{hcon} —la consistencia— es lo único que se pone de fuera.

\begin{leccion}[Por qué el teorema está partido en dos declaraciones]
Nótese que este enunciado no habla de \emph{la} sentencia de Gödel: habla de \textbf{cualquier}
punto fijo $G$. Es el argumento, aislado de la construcción que lo alimenta.

Eso no es elegancia gratuita. Cuando la construcción del punto fijo resultó estar apoyada en un
axioma inconsistente y hubo que rehacerla entera (capítulo~\ref{cap:numerales}), \textbf{este
teorema no se tocó}: bastó con enchufarle el punto fijo nuevo. Separar el argumento de su insumo
es lo que hizo que una reparación de los cimientos no obligara a rehacer el tejado.
\end{leccion}

Y la descarga, que es una línea:

\leanfrag{goedel_first_numeral}

\begin{observacion}
\emph{Cero postulados gödelianos.} No hay ninguna hipótesis del tipo «supongamos que $\Prov$
representa fielmente la demostrabilidad», que es donde las formalizaciones ingenuas esconden el
trabajo. Lo único que entra es la consistencia. Todo lo demás son teoremas de las partes
anteriores.
\end{observacion}

\section{Qué consistencia, exactamente}

La hipótesis se llama \ident{ConsistentOmega}, y conviene mirarla porque su nombre engaña:

\leanfrag{ConsistentOmega}

\selee{«la teoría no demuestra $\bot$ en el cálculo con la $\omega$-regla»}

Es \textbf{consistencia simple} —la más débil que se puede pedir—, enunciada sobre el cálculo
fuerte $\derives$. Y de ahí se hereda para el cálculo finitario:

\leanfrag{consistentH_of_omega}

\begin{muro}[Dos nombres a una palabra de distancia, y no son lo mismo]
En este proyecto conviven \ident{ConsistentOmega} y \ident{OmegaConsistent}. Se leen casi igual y
son propiedades distintas — y la que \emph{suena} más fuerte es la más débil:

\leanfrag{OmegaConsistent}

\ident{ConsistentOmega} es «el cálculo $\omega$ es consistente»: no demuestra $\bot$.
\ident{OmegaConsistent} es la \textbf{$\omega$-consistencia} de verdad: que la teoría no demuestre
$\exists x.\,A(x)$ mientras refuta $A(\bar n)$ para todos los testigos estándar.

La primera es la hipótesis de este capítulo. La segunda hace falta para la otra mitad. Confundirlas
en una lectura rápida convierte un teorema honesto en una exageración, y el orden de las palabras
es lo único que las separa.
\end{muro}

\section{Sobre qué teoría se ha demostrado esto}

\avisoreparacion

Conviene ser preciso sobre qué significa eso aquí. \ident{goedel_first_numeral} es un teorema
\textbf{real} sobre la teoría \textbf{reparada}: se retiró el axioma que hacía derivable $\bot$, y
lo que queda no tiene ninguna inconsistencia conocida. Pero «ninguna conocida» no es «ninguna», y
este libro no va a decir lo segundo.

\begin{leccion}[Y la hipótesis se lleva bien con esa cautela]
Hay una elegancia real en que el teorema tome la consistencia como \emph{hipótesis} en vez de
darla por buena. El enunciado no afirma que la teoría sea consistente: afirma que
\textbf{si} lo es, $G$ no es demostrable. Esa forma condicional es la que sobrevive a un
descubrimiento como el de julio de 2026 — y, de hecho, es la única forma en la que el segundo
teorema de Gödel permite enunciar nada sobre la consistencia de un sistema desde dentro.
\end{leccion}

\section{La mitad que falta}

\avisomediagodel

El teorema clásico dice que $G$ es \textbf{indecidible}: ni demostrable ni refutable. Lo que este
proyecto tiene demostrado es la primera mitad. La segunda existe como enunciado y arrastra una
hipótesis que hoy no está descargada — es el asunto del capítulo siguiente.

La razón de fondo no es un descuido, y ya estaba en el teorema original.

\begin{leccion}[Las dos mitades nunca costaron lo mismo]
Gödel demostró $\noderives G$ con la consistencia simple, y necesitó la
\textbf{$\omega$-consistencia} para $\noderives \lnot G$. La asimetría es del argumento, no de la
formalización: la primera mitad sólo tiene que impedir que la teoría demuestre algo y su negación;
la segunda tiene que impedir que la teoría se equivoque \emph{sobre los testigos}, y eso es una
exigencia más fuerte.

Rosser la eliminó en 1936 —cambiando la sentencia, no el argumento: su $R$ dice «si hay una
demostración mía, hay antes una demostración de mi negación»— y consiguió las dos mitades con
consistencia simple.\footnote{J.~B. Rosser, «Extensions of some theorems of Gödel and Church»,
\emph{The Journal of Symbolic Logic} 1(3), 1936, pp.~87--91.} Este proyecto \textbf{no} tomó esa
ruta: construyó la sentencia de Gödel y dejó la segunda mitad como una hipótesis explícita. Es una
decisión, y decirla es parte de decir qué se ha demostrado.
\end{leccion}

\section{Lo que queda dicho}

\begin{center}\small
\begin{tabularx}{\linewidth}{@{}l >{\raggedright\arraybackslash}X l@{}}
\toprule
mitad & qué hace falta & estado \\
\midrule
$\noderives G$ & consistencia simple, y nada más & \textbf{demostrado} \\
$\noderives \lnot G$ & reflexión —o $\omega$-consistencia más un verificador negativo— & hipótesis sin descargar \\
\bottomrule
\end{tabularx}
\end{center}

Media indecidibilidad no es un resultado a medias: $\noderives G$ es, por sí solo, el enunciado que
hunde el programa de Hilbert. Hay una sentencia aritmética que la teoría no puede demostrar, y su
número de Gödel se puede calcular. Lo que falta es saber si la teoría tampoco puede refutarla — y
en eso, como se verá, el proyecto tiene el enunciado escrito y ni una línea de la construcción.
                  % 13
% \chapter{La mitad que falta}
\label{cap:mitad-que-falta}

El capítulo anterior demostró que $G$ no tiene demostración. Para la indecidibilidad falta la otra
mitad: que tampoco tiene refutación. Este capítulo la construye hasta donde llega —que es lejos— y
después hace algo que un libro normal evitaría: explicar con precisión qué es lo que falta, por qué
no se puede coger el atajo, y cuánto trabajo hay hecho.

Porque el interés de esta mitad no es que esté incompleta. Es \emph{dónde} se atasca.

\section{El argumento, hasta donde llega}

Supóngase que $\lnot G$ sí fuera demostrable.

\begin{matematica}[Los cuatro pasos, y el quinto]
\begin{enumerate}
\item De $\Prf(\lnot G)$ se pasa a $\text{\ident{axioms}} \derives \lnot G$ por el puente.
\item El punto fijo, leído de derecha a izquierda, dice
  $\lnot\Prov(\codigo{G}) \Rightarrow G$. Con (1), la teoría demuestra
  $\lnot\lnot\Prov(\codigo{G})$.
\item Por eliminación de la doble negación, la teoría demuestra $\Prov(\codigo{G})$.
\item Y ahora hace falta el paso que no es lógica: de que \emph{la teoría demuestre que $G$ es
  demostrable} concluir que $G$ \textbf{realmente} lo es. Con eso, contradicción con el capítulo
  anterior.
\end{enumerate}
\end{matematica}

El paso (4) es todo el problema. Los tres primeros son manipulación formal; el cuarto dice que la
teoría no se equivoca al afirmar que algo es demostrable. Y eso no se sigue de que sea consistente.

\section{La hipótesis, con nombre}

El proyecto no la esconde: la declara.

\leanfrag{Reflects}

\selee{«\ident{Reflects} $G$ dice: si la teoría demuestra que $G$ es demostrable, entonces $G$ tiene
demostración»}

Es la \textbf{representabilidad negativa}, el espejo de D1. D1 dice que si algo es demostrable, la
teoría lo sabe; la reflexión dice que si la teoría lo afirma, es verdad. La primera es un teorema
del capítulo~\ref{cap:d1}. La segunda no lo es, y no puede serlo.

Con ella delante, la mitad se cierra:

\leanfrag{goedel_first_unrefutable}

\leanfrag{goedel_first_undecidable}

\selee{«bajo consistencia y reflexión, $G$ no es demostrable ni refutable»}

\section{Por qué el atajo está cerrado — y quien lo cierra es Gödel}

La tentación evidente es demostrar la reflexión dentro de la teoría. Haría falta que, para cada
$\varphi$ indemostrable, la teoría probara $\lnot\Prov(\codigo{\varphi})$. Los dos casos que uno
probaría primero son los dos que lo cierran:

\begin{muro}[Los dos callejones, y son el mismo]
Tómese $\varphi = \bot$. Si la teoría es consistente, $\bot$ es indemostrable, así que haría falta
$\derives \lnot\Prov(\codigo{\bot})$. Pero esa fórmula \textbf{es} $\mathrm{Con}(T)$, la sentencia
de consistencia — y por el segundo teorema de Gödel la teoría no la demuestra.

Tómese $\varphi = G$. Por el punto fijo, $\derives \lnot\Prov(\codigo{G})$ equivale a
$\derives G$ — y el capítulo anterior acaba de demostrar que eso no pasa.

La vía «que la teoría refute la demostrabilidad» no está cerrada por falta de trabajo:
\textbf{está cerrada por los teoremas de este libro}.
\end{muro}

\begin{leccion}[Un obstáculo que se demuestra es información, no un hueco]
Merece la pena parar en la forma de ese argumento. No dice «no lo hemos conseguido»: dice
\emph{dónde} está el obstáculo y \emph{qué} teorema lo pone. Un lector puede comprobarlo. Y como
efecto lateral, explica por qué existe Rosser: para tener las dos mitades desde consistencia simple
hay que cambiar de sentencia, porque con \emph{esta} sentencia el camino corto lo corta Gödel~II.
\end{leccion}

\section{El postulado que escondía todo esto}

Esta mitad estuvo dada por cerrada durante tres meses, y no lo estaba.

\begin{refutado}[Lo que decía la capa antigua]
El 13 de junio de 2026 la mitad $\noderives\lnot G$ se cerró apoyándose en un postulado escrito
como \textbf{bicondicional}:
\[ \bigl(\text{\ident{axioms}} \derives \Prov(\codigo{\varphi})\bigr)
   \;\Longleftrightarrow\;
   \bigl(\text{\ident{axioms}} \derives \varphi\bigr). \]
La dirección de derecha a izquierda es D1, y es un teorema. La de izquierda a derecha \textbf{es la
reflexión}, y es falsa en general. Y el teorema que se apoyaba en ella estaba enunciado bajo
\emph{consistencia simple}: es decir, afirmaba más de lo que Gödel permite. Se retiró.
\end{refutado}

\begin{leccion}[El mismo modo de fallo, por tercera vez]
Un enunciado plausible, aceptado como postulado, que resulta ser más fuerte de lo que se puede
sostener. Ya pasó con un axioma del kernel (capítulo~\ref{cap:kernel}) y con el axioma que hacía
inconsistente la teoría (capítulo~\ref{cap:numerales}). Aquí no produjo una contradicción: produjo
algo más difícil de detectar, \textbf{un teorema que decía de más}. El compilador no protesta
cuando una hipótesis es demasiado fuerte — sólo cuando falta.

La reparación no fue demostrar el postulado. Fue \textbf{partirlo}: quedarse con la mitad que es
teorema y dejar la otra a la vista como hipótesis.
\end{leccion}

\section{Cómo se descarga: dos piezas}

La reflexión no es indemostrable en absoluto — es indemostrable \emph{dentro} de la teoría. Desde
la metateoría se descarga, y el proyecto tiene el camino escrito:

\leanfrag{reflects_of_omega}

Hacen falta dos cosas. La primera es la \textbf{$\omega$-consistencia} —la de verdad, la que el
capítulo anterior distinguía del nombre parecido—. La segunda es un verificador que además de
aceptar las derivaciones buenas \textbf{rechace} las malas:

\leanfrag{NegVerifier}

\selee{«para toda $\varphi$ no demostrable y toda cadena estándar, la teoría refuta que esa cadena
verifique $\varphi$»}

Es el espejo negativo de D1: D1 dice que el verificador acepta lo que debe; esto dice que rechaza
lo que debe. Con las dos, la indecidibilidad sale sin hipótesis meta:

\leanfrag{goedel_first_undecidable_omega}

\begin{observacion}
La demostración de \ident{reflects_of_omega} es donde vive el \textbf{único uso esencial de lógica
clásica} de todo el proyecto, y el capítulo~\ref{cap:objeto-y-meta} lo identificó: es el principio
de Markov. Tiene sentido que esté aquí. Lo que se hace es concluir «existe una demostración» de la
imposibilidad de que no exista, y eso —sobre un enunciado $\Sigma_1$ con matriz decidible— es
exactamente Markov.
\end{observacion}

\section{Y ahora la cifra incómoda}

\ident{NegVerifier} está escrito. Está enunciado con precisión, tiene su papel claro en la cadena y
dos teoremas lo consumen como hipótesis.

\textbf{Y no había ni un teorema en producción que lo concluyera.} Así se midió al escribir este
capítulo, la mañana del 10 de septiembre de 2026: cero.

\begin{observacion}[Corregido esa misma tarde, y las dos mitades de la corrección importan]
La frase de arriba se quedó a medias en unas horas, y la parte que se cayó no es la que parece.

Lo que \textbf{sigue siendo cierto} es que ningún teorema concluye \ident{NegVerifier} de manera
incondicional. Lo que \textbf{dejó de serlo}: ya no es cierto que no haya nada. Existe
\ident{negVerifier_of_deudas}, que lo concluye \textbf{a partir de dos obligaciones enunciadas},
\ident{DEUDA_chainNeg} y \ident{DEUDA_inNeg}; y la solidez estructural del verificador —el módulo
que el plan llamaba «el corazón» y estimaba en 300--500 líneas— resultó estar \textbf{ya
demostrada}: \ident{verifier_sound} \emph{es} \ident{decodeChain_prf}.

Y lo que era \textbf{directamente falso} es la causa que esta sección daba: «lo que lo bloquea
—la inyectividad de la codificación numérica— vive en un sondeo, fuera del build». No: \ident{consN_inj}
y \ident{codeNat_ne} están en producción desde el 2 de septiembre. Lo dice la propia auditoría de
este libro, y es un buen ejemplo de por qué §2.6 vale también \emph{hacia dentro}: la frase se
copió de un plan del proyecto en vez de medirse.

El bloqueo real era otro, y hubo que medirlo: la clase de testigos admitía basura cuya refutación
exigía \textbf{evaluar el emparejamiento de Cantor}, que es inviable. Se estrechó
(\textsc{adr}-022), con el precio escrito —la $\omega$-consistencia queda algo más fuerte— y con
la garantía que lo hace admisible: el código de \emph{cualquier} derivación real sigue siendo un
testigo de la clase.
\end{observacion}

\avisomediagodel

\begin{leccion}[Qué es exactamente lo que falta, y por qué decirlo así vale más]
Compárese con la manera habitual de dejar esto: «la otra mitad es análoga» o, peor, «Gödel~I,
completo». Lo que este proyecto puede decir en su lugar es:

\begin{center}\small
\begin{tabularx}{\linewidth}{@{}l >{\raggedright\arraybackslash}X@{}}
\toprule
pieza & estado \\
\midrule
el argumento $\noderives\lnot G$ desde la reflexión & \textbf{demostrado} \\
la reducción reflexión $\leftarrow$ $\omega$-consistencia $+$ verificador negativo & \textbf{demostrada} \\
el verificador negativo & \textbf{reducido a dos obligaciones con nombre} \\
\bottomrule
\end{tabularx}
\end{center}

Eso no es un hueco: es un plano con una pieza marcada en rojo. La diferencia entre las dos maneras
de escribirlo es la diferencia entre un libro que se puede auditar y uno que hay que creer.
\end{leccion}
       % 14
% \chapter{Las tres condiciones de derivabilidad}
\label{cap:d123}

El primer teorema necesita que $\Prov$ exista y diga la verdad en una dirección. El segundo
necesita mucho más: que la teoría sepa \emph{razonar} sobre su propia demostrabilidad. Esas
exigencias son tres, se llaman D1, D2 y D3, y no las publicó Gödel — son el artículo que anunció y
nunca escribió, y las suministraron Hilbert y Bernays en 1939.

Este capítulo dice qué pide cada una, cuáles son teoremas aquí y por qué la tercera lleva meses
resistiéndose. Y da una respuesta concreta a la pregunta que el lector debería estar haciéndose:
\emph{¿cuánto falta, exactamente?}

\section{Las tres, en una página}

\begin{matematica}[Lo que pide cada una]
\[
\begin{array}{ll}
\text{D1} & \vdash \varphi \;\Longrightarrow\; \derives \Prov(\codigo{\varphi}) \\[4pt]
\text{D2} & \derives \Prov(\codigo{\varphi \Rightarrow \psi}) \Rightarrow
             \bigl(\Prov(\codigo{\varphi}) \Rightarrow \Prov(\codigo{\psi})\bigr) \\[4pt]
\text{D3} & \derives \Prov(\codigo{\varphi}) \Rightarrow \Prov(\codigo{\Prov(\codigo{\varphi})})
\end{array}
\]

\noindent D1: \emph{si algo se demuestra, la teoría lo sabe.} \quad
D2: \emph{la demostrabilidad respeta el modus ponens.} \quad
D3: \emph{la teoría sabe que sabe.}
\end{matematica}

La diferencia de naturaleza entre la primera y las otras dos es la que gobierna todo el capítulo, y
conviene fijarla antes de mirar ningún código.

\begin{leccion}[Externo contra provable]
\textbf{D1 es una afirmación de la metateoría}: para \emph{cada} $\varphi$ demostrable,
\emph{nosotros} construimos la derivación de $\Prov(\codigo{\varphi})$. Se demuestra caso a caso, y
cada caso es finito.

\textbf{D2 y D3 son afirmaciones de la teoría objeto}: la que se demuestra es una fórmula con un
$\Rightarrow$ dentro, cuantificada sobre \emph{todas} las derivaciones posibles. Y razonar sobre
todas las derivaciones, dentro de la teoría, exige inducción — que es justo lo que
Q\texttt{++} no tiene.

Ésa es toda la historia. D1 es trabajo; D3 es un problema.
\end{leccion}

\section{D1: un teorema, sin hipótesis}

Ya está demostrada, y el capítulo~\ref{cap:d1} contó cómo. Aquí sólo el recordatorio de que no
arrastra nada:

\leanfrag{repr_pos_prf}

Es la $\Sigma_1$-completitud \textbf{externa}: dada una demostración concreta, se construye el
testigo concreto. Ninguna inducción de la teoría objeto interviene.

\section{D2: también un teorema, y es el que suele postularse}

\leanfrag{d2_prf}

\selee{«la teoría demuestra que, si es demostrable $A \Rightarrow B$ y es demostrable $A$, entonces
es demostrable $B$»}

\begin{observacion}
Merece la pena subrayarlo porque es donde muchas presentaciones —y una versión anterior de este
mismo proyecto— se ahorran el trabajo: \textbf{D2 aquí es un teorema, no un postulado}. La
diferencia se ve en el inventario: el capítulo~\ref{cap:lean} decía que un \ident{axiom} es
una deuda, y ésta está pagada.

La prueba, por cierto, es instructiva: elimina el $\exists$ externo del predicado de
demostrabilidad —el testigo es la propia derivación— y después concatena las dos cadenas. Es el
\emph{modus ponens} de la teoría, escrito dentro de la teoría.
\end{observacion}

\section{D3: la que Gödel anunció, y estuvo postulada hasta anteayer}

Durante casi toda la vida de este proyecto, D3 fue lo que la mayoría de las presentaciones hacen
con ella: un postulado. Estaba escrita como \ident{axiom} de Lean —no como axioma objeto de la
lista de 141, sino como un supuesto de la \textbf{metateoría}—, lo que significaba que todo
teorema que la citara la llevaba impresa en su \emph{footprint}. El segundo teorema de Gödel se
publicaba con esa marca.

Ya no. Desde el 10 de septiembre de 2026 es un teorema, con la misma firma exacta que tenía el
axioma:

\leanfrag{d3_teorema}

\selee{«la teoría demuestra que, si $\varphi$ es demostrable, entonces es demostrable que
$\varphi$ es demostrable»}

\begin{leccion}[Retirar un axioma sólo puede fortalecer lo que había]
Nótese qué es lo que \emph{no} cambió: ni el enunciado de D3, ni el de \ident{goedel_second'}, ni
una línea de su demostración. Lo único que cambió es de dónde viene D3 — antes se suponía, ahora se
deriva.

Ésa es la propiedad que hace que este tipo de deuda sea \textbf{segura de contraer}: mientras el
postulado esté a la vista y con la firma que tendrá el teorema, pagarlo después no obliga a
reconstruir nada aguas abajo. Se sustituye la palabra \ident{axiom} por \ident{theorem} y todos
los \emph{footprints} de la cadena adelgazan solos.
\end{leccion}

El inventario de axiomas de Lean del proyecto pasó, ese día, de siete a \textbf{seis}. De los que
quedan, ninguno es una deuda de este tipo: son anclas de codificación y esquemas de inducción
declarados.

\section{Por qué costó tanto}

La demostración real de D3 es la $\Sigma_1$-completitud \textbf{provable}: que la teoría
demuestre, de toda derivación, que su verificador la acepta. Y ahí hay que razonar sobre un objeto
de longitud arbitraria \emph{desde dentro}.

\begin{leccion}[Dónde estaba la dificultad, en una frase]
D1 mira \emph{una} derivación y construye \emph{un} testigo. D3 tiene que decir algo sobre
\textbf{todas} las derivaciones, y decirlo la teoría, no nosotros. La primera es una función; la
segunda es una inducción — y la inducción de la teoría objeto es exactamente el recurso que
Q\texttt{++} no tiene y que \capfuturo{la inducción como precio} tendrá que comprar.
\end{leccion}

\section{La forma de la demostración}

Lo que hay debajo del teorema de una línea es una reducción de varios pisos, y su forma final es
ésta: D3 se sigue de reflejar \textbf{las dos mitades del cuerpo} del verificador.

\leanfrag{d3_prf_of_halves}

\selee{«si se reflejan las dos mitades del cuerpo —la buena formación de cada línea y sus
premisas—, D3 queda demostrada»}

Las dos reciben lo mismo: una cadena $q$ y un índice $i$ dentro de ella. Y de la cadena sale lo que
las dos necesitan:

\leanfrag{prf_lineWF_of_chainOk}

\begin{center}\small
\begin{tabularx}{\linewidth}{@{}l >{\raggedright\arraybackslash}X l@{}}
\toprule
mitad & qué refleja & se cerró \\
\midrule
(a) & que la línea $i$-ésima está bien formada & 2026-09-10 \\
(b) & que sus premisas están donde deben & 2026-09-10 \\
\bottomrule
\end{tabularx}
\end{center}

La mitad (a) es la que costó el episodio del capítulo~\ref{cap:muro-substfc}: los veintiún
reflectores del verificador, de los que siete estuvieron bloqueados cuarenta días. La mitad (b)
necesitó tres piezas más, y la primera de ellas resume bien de qué va todo esto:

\leanfrag{pcc_eval_premsOf}

\selee{«si la línea $t$ está bien formada, la teoría demuestra que el operador \emph{premisas de}
aplicado a su código da el código de sus premisas»}

Veintiuna ramas, una por cada tipo de línea que el verificador admite, porque \textbf{las premisas
sólo se pueden evaluar sabiendo de qué regla se trata}. Con ella, la cota; con la cota, el chasis
inductivo interior; y con las tres, el ensamblaje:

\leanfrag{d3_prf_real}

\begin{observacion}
Este libro llegó tarde por dos días. El capítulo que se escribió el 10 de septiembre por la mañana
decía —midiendo, y correctamente— que la mitad (b) estaba abierta en tres piezas, y señalaba además
un hueco de ensamblaje que ningún documento del proyecto había visto: la mitad (a) estaba probada y
\textbf{no la consumía nadie}.

Las tres piezas se cerraron esa misma tarde, en cuatro commits, y el hueco de ensamblaje lo cerró
el último de ellos. \ident{d3_prf_real} es exactamente el teorema que faltaba: la mitad (a) y la
mitad (b), enchufadas.

Se deja escrito porque es la clase de dato que un libro sobre un proyecto vivo debe dar y casi
ninguno da: \textbf{cuánto tiempo hubo entre saber qué faltaba y que dejara de faltar}. Aquí,
unas horas.
\end{observacion}

\section{Lo que las tres condiciones significan juntas}

\begin{leccion}[Por qué esto es el artículo que Gödel no escribió]
El primer teorema necesita D1 y nada más. El segundo necesita las tres, y ésa es la razón de que
Gödel anunciara una segunda parte y no la publicara: la demostración detallada exige justamente
esto, y «esto» resultó ser considerablemente más trabajo de lo que cabía en una nota.

Noventa y cinco años después, en una formalización mecánica, la proporción se repite con precisión
incómoda: D1 y D2 salieron en semanas, y D3 —la que Gödel dejó anunciada— costó meses y fue la
última en caer. Eso no dice nada malo del proyecto: es la medida honesta de lo que Hilbert y
Bernays tuvieron que suministrar cuando la segunda parte no llegó.

Con las tres demostradas, el segundo teorema deja de publicarse con una condición al margen. Qué
significa eso exactamente —y qué hipótesis \emph{sí} le quedan, que no son ninguna de las
tres— es el asunto de \capfuturo{Gödel II}.
\end{leccion}
% 15
% \include{capitulos/cap-la-induccion-como-precio} % 16
% \include{capitulos/cap-godel-ii-modulo-d3}       % 17

% =============================================================================
\part{Lo que no sale en los libros}
\label{parte:noescrito}

\noindent \emph{¿Y qué pasó de verdad mientras se construía todo esto?} Una inconsistencia que sobrevivió meses, cómo se localizó, cuatro reparaciones que no funcionan, la que sí, y lo que se aprendió. Es la parte que ningún manual escribe, porque los manuales presentan el resultado ya limpio.
% =============================================================================
\chapter{El muro de \texorpdfstring{\ident{substfc}}{substfc}}
\label{cap:muro-substfc}

Los libros de incompletitud cuentan la demostración. Esta parte cuenta lo otro: qué pasa cuando
la pieza que falta no es difícil, sino que \textbf{la ruta obvia hacia ella hace inconsistente la
teoría}, y hay que descubrirlo antes de escribirla.

El episodio de este capítulo duró cuarenta días. No fue un atasco de implementación: fue un muro
con tres paredes independientes, y la primera de ellas —la que parecía la salida— habría destruido
el proyecto entero de haberse aceptado. Se atrapó porque alguien se sentó a derivar la
contradicción a mano en vez de comprobar que compilaba.

\section{Qué es \texorpdfstring{\ident{substfc}}{substfc}, y por qué está en todas partes}

El capítulo~\ref{cap:codigos} construyó los códigos de la sintaxis. Sustituir es lo que
la sintaxis hace: si $\varphi$ tiene una variable libre, $\varphi[t/v]$ es la fórmula que resulta de
poner $t$ donde estaba $v$. En la metateoría eso es la función \ident{substFormula}. Dentro de la
teoría objeto tiene que existir su contraparte, y se llama \ident{substfc}:

\leanfrag{substfc_def}

\selee{«\ident{substfc} $v$ $t$ $f$ es el término objeto \emph{«sustituir la variable de código $v$
por el término de código $t$ en la fórmula de código $f$»}»}

Nótese qué es y qué no es. No computa nada: construye un \emph{símbolo de función objeto} de aridad
tres. Lean no sabe qué significa. Quien lo sabe es la teoría, mediante ocho axiomas objeto, uno por
constructor de fórmula. Éste es el del condicional:

\leanfrag{ax_substfc_impl}

\selee{«para cualesquiera $v,t,a,b$: sustituir en el condicional $a \Rightarrow b$ es el condicional
de las dos sustituciones»}

Y aquí está la razón de que todo esto importe. El término diagonal —el corazón mismo de la sentencia
de Gödel, la que dice de sí misma que no es demostrable— \emph{es} una aplicación de
\ident{substfc}:

\leanfrag{diagTerm}

Si \ident{substfc} no funciona, no hay autorreferencia. Si tiene una grieta, la grieta está en el
centro del edificio.

\section{Lo que bloqueaba: siete de veintiuno}

D3 —la tercera condición de derivabilidad, la que Gödel anunció y nunca publicó, y sin la cual no
hay segundo teorema— se reduce en este proyecto a un solo lema, y ése a que el verificador de
pruebas se refleje dentro de $\Prov$: que la teoría demuestre, de cada línea de una derivación, que
la línea está bien formada.

Las líneas se clasifican por un \emph{tag}, y hay veintiún tags. Cada uno necesita su
\defterm{reflector}\nivelterm{proyecto}: el teorema que traslada al interior de $\Prov$ lo que la
metateoría ya sabe de ese tag. Catorce estaban hechos. Siete no.

Y aquí el proyecto hizo algo que conviene subrayar, porque es lo contrario de lo que se suele
hacer con una tarea pendiente: en vez de estimar cuánto faltaba, lo \textbf{demostró}.

\leanfrag{lineWF_modulo_7}

\begin{leccion}[Un teorema cuya función es medir lo que falta]
Ese enunciado no prueba nada nuevo sobre aritmética. Prueba que \emph{el ensamblaje encaja}: que
tomando los siete reflectores que faltan \textbf{como hipótesis}, todo lo demás cierra. Y por tanto
que cerrar el verificador es \textbf{exactamente} cerrar esos siete, ni una pieza más aguas abajo.

Es una manera de convertir «creo que falta poco» en un objeto que compila o no compila. Cuesta una
tarde y evita meses de trabajo en la dirección equivocada — que es justo lo que este capítulo
narra que estuvo a punto de pasar.
\end{leccion}

Los siete —\ident{q1}, \ident{q2}, \ident{q3}, \ident{leibniz}, \ident{ind}, \ident{qconf},
\ident{listInd}— tienen algo en común y sólo eso: su conclusión se reconstruye con \ident{substfc}
o con \ident{liftfc}, y por tanto exige \emph{evaluar} esos operadores dentro de $\Prov$. Eso es el
muro.

\section{El axioma que lo resolvía, y que hace inconsistente la teoría}

Hay una salida evidente, y es la primera que se le ocurre a cualquiera. El operador dotado
\ident{substfcT} es la versión «código del código» de \ident{substfc}. Bastaría postular que
conmutan:
\[
  \text{\ident{ax_tc_substfc}}:\quad
  \dotado{(\text{\ident{substfc}}\ v\ s\ f)} \;\doteq\;
  \text{\ident{substfcT}}\ \dotado{v}\ \dotado{s}\ \dotado{f}.
\]

\selee{«el código del resultado de sustituir es el resultado de sustituir sobre los códigos»}

Parece inocente. Ya existe un axioma exactamente de esa forma y nadie protesta — el de \ident{cons}:

\leanfrag{ax_tc_cons}

La diferencia es la que decide el episodio, y no se ve mirando los dos enunciados: \ident{cons} es
un \textbf{constructor}, y \ident{substfc} es una \textbf{función definida} — definida, además, por
ocho ecuaciones que la teoría \emph{ya tiene} y que la interpretan como productora de árboles de
\ident{cons}. Postular la ley de conmutación sobre algo ya interpretado no añade información:
la contradice.

\begin{refutado}[La derivación, en cinco pasos]
Supóngase \ident{ax_tc_substfc} en la teoría. Tómese la fórmula $\text{\ident{implc}}\ a\ b$.
\begin{enumerate}
\item Por \ident{ax_substfc_impl} —axioma ya presente—, la teoría demuestra que
  $\text{\ident{substfc}}\ v\ s\ (\text{\ident{implc}}\ a\ b)$ \emph{es} un
  $\text{\ident{implc}}$ de dos sustituciones. La teoría \textbf{interpreta} \ident{substfc}: no
  es opaco.
\item El operador $\dotado{\cdot}$ respeta la igualdad derivada —es un símbolo de función, y la
  congruencia de Leibniz vale para él—:
  \leanfrag{prf_congr_tcFn}
  Luego los códigos de los dos lados de (1) son iguales.
\item El lado izquierdo, por el axioma supuesto, tiene cabeza
  $\langle 1, \codigo{\text{\ident{substfc}}}, \ldots\rangle$.
\item El lado derecho es el código de un \ident{cons} —porque \ident{implc} \emph{es} un
  \ident{cons}—, y por \ident{ax_tc_cons} tiene cabeza
  $\langle 1, \codigo{\ident{::}}, \ldots\rangle$.
\item Luego la teoría demuestra que esas dos cabezas son iguales. Pero la aritmética negativa de
  códigos, que ya está construida, demuestra que \textbf{no} lo son:
  \leanfrag{cons_ne_head}
  Contradicción. La teoría prueba $\bot$.
\end{enumerate}
\end{refutado}

\begin{muro}[Por qué esto no habría fallado en compilación]
Una teoría objeto inconsistente \textbf{compila}. Lean no comprueba que los axiomas objeto sean
satisfacibles: son datos, valores de tipo \ident{Formula}, y añadir uno más a la lista es tan legal
como añadir el primero. El proyecto habría seguido en verde durante meses, demostrando teoremas
—todos ellos vacíos, porque de una teoría inconsistente se demuestra cualquier cosa—. Y en
particular habría «demostrado» $\Prov(\codigo{\bot})$, que aniquila la construcción de Gödel entera:
la sentencia $G$ dejaría de decir nada.

Este es el modo de fallo característico del proyecto, y por eso la Parte~\ref{parte:noescrito}
existe. Ya apareció en pequeño en el capítulo~\ref{cap:kernel} —un axioma del kernel que era falso—
y volverá a aparecer en grande en el capítulo~\ref{cap:numerales}, donde no se atrapó a tiempo.
\end{muro}

\begin{observacion}
\ident{ax_tc_substfc} \textbf{no existe} en el código. Y sin embargo su nombre aparece
\textbf{nueve veces} en el árbol —una en \modulo{ROBINSON_PlusPlus/Meta/EvalArithPrf.lean}, otra en
\modulo{ROBINSON_PlusPlus/Meta/SubstfcCodePrf.lean}, el resto en sondeos y experimentos—: siempre
en un comentario, siempre para advertir de que sería inconsistente. El proyecto convirtió la
refutación en un cartel y lo clavó en las nueve puertas por las que alguien podría volver a entrar.
Es, en pequeño, la misma doctrina que este libro: un resultado negativo también es un resultado, y
se guarda donde vaya a leerse.
\end{observacion}

\section{El obstáculo real: evaluar sobre un argumento abstracto}

Cerrada la puerta del axioma, queda el trabajo honesto: demostrar la evaluación en vez de
postularla. Y ahí está la segunda pared.

Hay evaluaciones provables que salen. La de \ident{carc} —la cabeza de una lista— sale porque
\ident{carc} tiene \textbf{una} ecuación definitoria y se instancia sobre un patrón explícito:
delante hay un \ident{cons} escrito, y se aplica la ecuación. \ident{substfc} no está en esa
situación por dos razones a la vez: se define por recursión sobre los \textbf{ocho} constructores
de fórmula, con ligadores —los casos $\forall$ y $\exists$ incrementan el nivel y desplazan el
sustituyendo—, y en los siete tags se aplica a un código \textbf{abstracto}, una casilla de la
línea que se está verificando, no a una fórmula escrita.

\begin{muro}[Lo que eso obliga a hacer]
Evaluar \ident{substfc} sobre un argumento abstracto exige \textbf{inducción sobre códigos de
fórmula dentro de $\Prov$}. No inducción en la metateoría —eso se tiene— sino inducción de la
teoría objeto sobre sus propios códigos, disponible en un punto en que la teoría objeto no tiene
inducción: la aritmética de \modulo{ROBINSON_PlusPlus/Minimal/} es Q\texttt{++}, sin esquema de
inducción, por la razón que dio el capítulo~\ref{cap:qpp}.

Ésa es la distancia real entre «un lema más» y «un problema abierto», y explica por qué el registro
del proyecto insiste en que el muro no era mecánico.
\end{muro}

\section{La tercera pared: quitar \texorpdfstring{\ident{substfc}}{substfc} tampoco se puede}

Queda una salida más, y es la que parece más razonable de todas: si el problema es que los siete
esquemas mencionan \ident{substfc}, reescríbanse sin él. El 25 de julio de 2026 se sondeó
sistemáticamente, y el resultado no fue «no lo hemos conseguido» sino una dicotomía cerrada.

El argumento tiene dos bisagras, las dos comprobables en el código: la regla de instanciación toma
un término \textbf{arbitrario}, y \ident{substFormula} es recursión sobre los ocho constructores.
De ahí: para que el reflector sea \emph{sólido} tiene que atar la instancia correcta, y la instancia
correcta contiene $\varphi[t/v]$. Si esa mención queda \textbf{ligada} —como función, como ecuación,
como átomo relacional—, reaparece el mismo muro un piso más abajo. Si queda \textbf{libre}, el
esquema deja de ser sólido. Las tres propuestas concretas que se estudiaron colapsan las tres al
mismo núcleo.

\begin{leccion}[La diferencia entre «no sale» y «no existe»]
Un resultado como éste es incómodo de publicar y valiosísimo de tener: no dice que el equipo no
supiera hacerlo, dice que no hay nada que hacer por ahí. Sin él, la ruta se habría reintentado cada
pocas semanas con una variante — que es exactamente lo que pasa cuando un callejón sin salida no se
documenta como tal. Los otros casos de esta misma clase que el proyecto encontró después los
recoge \capfuturo{las tareas que resultan ser imposibles}.
\end{leccion}

\section{Cuarenta días}

\begin{center}\small
\begin{tabular}{lll}
\toprule
fecha & commit & hito \\
\midrule
2026-07-10 & \texttt{8a96e8e} & el frente arranca en positivo \\
\textbf{2026-07-22} & \texttt{15dccda} & \textbf{el muro se identifica}: «el plan era erróneo en su forma» \\
2026-07-24 & --- & veredicto: \ident{ax_tc_substfc} es inconsistente; no hay ruta \defterm{net-0}\nivelterm{proyecto} \\
2026-07-25 & \texttt{6c4d7ac} & la puerta de reescribir los esquemas, cerrada \\
2026-07-26 & \texttt{ee31d07} & inducción fuerte objeto, net-0 — la herramienta que luego serviría \\
2026-08-25 & \texttt{bdfff12} & el reconocedor fusionado no desbloquea el muro \\
2026-08-30 & \texttt{2f27a29} & el descenso, cerrado \\
\textbf{2026-08-31} & \texttt{c1fd804} & \textbf{el muro está roto} \\
2026-09-07 & \texttt{61edd46} & dentro del build: deja de ser un sondeo y pasa a ser un teorema \\
\bottomrule
\end{tabular}
\end{center}

Cuarenta días entre la identificación y la rotura; siete más hasta que el resultado dejó de vivir
en \modulo{sondeos/} y entró en el build, que es la condición que este libro exige para llamarlo
teorema (§2.2). Hoy se enuncia así:

\leanfrag{pcc_eval_substfc}

\selee{«bajo dos guardas —que $s$ tenga testigo y que $f$ sea un código de fórmula bien formado
contra su testigo—, la teoría demuestra que el operador dotado aplicado a los códigos es igual al
código del resultado»}

Las dos hipótesis no son deudas colgando: son las \emph{guardas} que dicen sobre qué códigos vale
la evaluación, y \capfuturo{la no-vacuidad de las guardas} muestra que se satisfacen
para códigos reales. Cómo se llegó hasta aquí —definir en vez de axiomatizar, partir el
reconocedor en tres, y una inducción con la guarda dentro— es la materia del resto de esta
parte.

\begin{observacion}
Un detalle que este libro no puede pasar por alto, porque es su propia doctrina aplicada a sí mismo.
El docstring de \ident{pcc_eval_substfc} dice, literalmente, «cero \ident{sorry}, cero hipótesis».
El teorema tiene dos. Es una imprecisión menor —«hipótesis» quiere decir ahí «obligación
pendiente»—, pero es exactamente el tipo de afirmación que un lector tomaría por buena y que el
compilador no verifica. Queda anotada en \doc{DOCSTRINGS-NO-FIABLES.md}.
\end{observacion}
               % 18 ✅
% \include{capitulos/cap-inconsistencia-latente}   % 19
% \include{capitulos/cap-localizar-el-dano}        % 20
% \include{capitulos/cap-cuatro-reparaciones}      % 21
\chapter{La reparación: códigos como numerales}
\label{cap:numerales}

El capítulo anterior terminó en un sitio incómodo: cuatro reparaciones razonables, y ninguna
funciona. Éste cuenta la que sí, y por qué —vista desde el final— era la única que podía funcionar.

\section{El diagnóstico, en una línea}

La inconsistencia no venía de un axioma equivocado sino de un \textbf{error de categoría}. El
símbolo objeto \ident{tcFn} —«el código del código»— es una operación sobre \emph{sintaxis},
declarada como función \emph{objeto}. Y una función objeto sólo puede depender de \textbf{valores}.

\begin{leccion}[La distinción que lo explica todo]
Un predicado \emph{es} un subconjunto: separa números, no maneras de escribirlos.
\ident{substfc} necesita un reconocedor \textbf{extensional} —¿pertenece este número al conjunto de
los códigos de fórmula?—, y eso un predicado lo puede dar. \ident{tcFn} necesita una distinción
\textbf{intensional} —¿qué sintaxis escribimos para denotar el número 9?—, y eso ningún predicado
lo puede dar, porque no está separado en los valores.
\end{leccion}

Con $\codigo{\varphi}$ escrito como \textbf{árbol} \ident{cons}, la sintaxis no se recupera del
valor: \ident{ax_L0_cons_def} identifica \ident{cons h t} con \ident{pair h (σt)}, así que en
$\mathbb{N}$ el mismo número es a la vez un par y un sucesor, y \ident{tcFn} debía devolver dos
códigos distintos para él. Con $\codigo{\varphi}$ escrito como \textbf{numeral}, la ambigüedad
desaparece por una razón que conviene decir despacio: \textbf{el numeral es canónico}. Sólo hay una
manera de escribir $\sigma^n 0$.

\begin{refutado}[Y no era «indemostrable»: era falso]
La lectura sintáctica de \ident{tcFn} sobre \ident{cons} con argumentos abstractos no es una ecuación
que faltase probar. Bajo la lectura numeral obliga a que un código de cabeza $\codigo{\sigma}$ sea
igual a uno de cabeza $\codigo{::}$. Medido en \modulo{sondeos/PilotoRastreada.lean}.
\end{refutado}

\section{La aritmética sale sin división}

El obstáculo práctico de la vía numeral es el polinomio de Cantor: \ident{cons} se define como
$\mathrm{div2}$ de él, y razonar sobre divisibilidad dentro de la teoría objeto es caro. La salida
es no razonar sobre ella en absoluto. Se define el valor del código con \textbf{números
triangulares}, de modo que la mitad ya está tomada.

\begin{matematica}[Números triangulares y valor de un \ident{cons}]
\[
T(0)=0,\qquad T(n{+}1)=T(n)+(n{+}1),
\qquad\text{luego}\qquad 2\,T(n)=n(n{+}1).
\]
\[
\mathrm{consN}(a,b) \;=\; T\bigl(a+(b{+}1)\bigr) + (b{+}1),
\qquad
2\cdot\mathrm{consN}(a,b) \;=\; \underbrace{(a{+}b{+}1)(a{+}b{+}2) + 2(b{+}1)}_{\text{polinomio de Cantor de }\langle a,b\rangle}.
\]
La segunda identidad es una igualdad de \ident{Nat}, demostrable por reescritura. No hay división
en ninguna parte.
\end{matematica}

\leanfrag{triN}
\leanfrag{consN}
\leanfrag{two_mul_consN}

\begin{leccion}[Dónde está el truco]
No se ha evitado la división eligiendo mejor las tácticas: se ha evitado \emph{eligiendo la
definición}. \ident{consN} se define ya multiplicado por dos, y por eso engancha directamente en la
forma $2m$ exacta que pide \ident{prf_div2_numeral}. Cuando una prueba obliga a razonar sobre una
operación cara, la primera pregunta no es cómo probarlo, sino si la definición se puede escribir
del otro lado.
\end{leccion}

\section{De la identidad en \ident{Nat} a la evaluación provable}

La identidad anterior vive en la metateoría. Para que sirva hay que subirla al cálculo: la teoría
objeto tiene que \emph{probar} que el árbol de código y el numeral de su valor son el mismo número.

\leanfrag{prf_cons_eval}

Con ese peldaño, el espejo meta de \ident{formCode} —la función \ident{codeNat}, que calcula en
\ident{Nat} lo que \ident{formCode} construye como término— cierra el puente por meta-recursión sobre
la estructura de la fórmula.

\leanfrag{codeNat}

\begin{teorema}[El código de una fórmula \emph{es} un numeral]
\label{teo:formcode-numeral}
Para toda fórmula $\varphi$, la teoría prueba
$\codigo{\varphi} \;=\; \mathrm{numeral}\bigl(\mathrm{codeNat}(\varphi)\bigr)$.
\end{teorema}

\leanfrag{prf_formCode_numeral}

Obsérvese que la prueba es una \textbf{meta-recursión}: un caso por constructor de \ident{Formula},
cada uno ensamblado con \ident{prf_cons_eval_of}. No hay inducción de la teoría objeto por ninguna
parte, y por eso el footprint se queda en la base del kernel.

\section{El puente, y qué cambia en los enunciados}

Del lado del cálculo $\derives$, el mismo hecho es la hipótesis que el piloto del lema diagonal
daba por supuesta, y que aquí queda \textbf{descargada}:

\leanfrag{hFN}

\begin{muro}[Lo que sí cambia]
Refundar por la vía numeral \textbf{cambia enunciados, no sólo pruebas}: el código estático que
aparece dentro de los teoremas pasa de \ident{termCode (formCode φ)} a
\ident{termCode (numeral (codeNat φ))}. Cada consumidor hay que revisarlo, o componerlo con
\ident{provCode_transfer}, que puentea las dos representaciones en un solo paso de Leibniz.
\end{muro}

\section{Cómo se decidió: pilotar antes de ejecutar}

Aquí está, para mi gusto, la lección de ingeniería del capítulo. La reparación exigía construir
unos veinte lemas nuevos. Antes de escribir ninguno se hizo lo contrario de lo intuitivo:
\textbf{se asumió el resultado costoso como axioma de Lean} y se comprobó que, con él, la cadena
cerraba de verdad.

\leanfrag{piloto_codeN_axioma}
\leanfrag{piloto_hFN_axioma}

Con esos dos axiomas puestos, el piloto reconstruyó el lema diagonal entero y comprobó que el punto
fijo salía. Sólo entonces se invirtió el trabajo en probarlos — y el sondeo gemelo descargó \ident{hFN}
sin axioma ninguno:

\leanfrag{descarga_hFN}

\begin{leccion}[Un axioma temporal es un instrumento de medida]
Poner un \ident{axiom} de Lean no es hacer trampa si se hace \emph{para preguntar}. Contesta la única
pregunta que importa antes de invertir: \emph{si tuviera esto, ¿cerraría?} Si la respuesta es no, se
han ahorrado veinte lemas; si es sí, se sabe exactamente dónde van a enchufarse. Lo que sería hacer
trampa es olvidarse de quitarlo — y contra eso está el inventario de \doc{AXIOMS.md} y la auditoría
por \printaxioms.
\end{leccion}

\section{El coste}

\begin{center}\small
\begin{tabular}{ll}
\toprule
Axiomas retirados & \ident{ax_tc_cons} \\
Axiomas nuevos & ninguno \\
Balance & $-1$ \\
Módulos afectados & 31, a cuarentena (capítulo siguiente) \\
\bottomrule
\end{tabular}
\end{center}

\noindent Es un resultado poco frecuente: la reparación de una inconsistencia que \textbf{reduce} la
superficie axiomática en vez de ampliarla. La razón es que no se añadió nada para tapar el agujero;
se cambió la \emph{representación} para que el agujero no existiera.

\avisoreparacion
     % 22 ✅
% \include{capitulos/cap-la-cuarentena}            % 23
% \include{capitulos/cap-la-reconstruccion}        % 24
% \include{capitulos/cap-definir-en-vez-de-axiomatizar} % 25
% \include{capitulos/cap-la-particion}             % 26
% \include{capitulos/cap-romper-el-muro}           % 27
% \include{capitulos/cap-tarea-imposible}          % 28
% \include{capitulos/cap-metodo}                   % 29

% =============================================================================
\appendix
% =============================================================================
% \include{capitulos/ap-inventario-de-axiomas}
% \include{capitulos/ap-mapa-de-modulos}
% \include{capitulos/ap-trampas-de-lean}
%% Apéndice: glosario. Los términos se INTRODUCEN narrativamente en la Parte I;
%% esto es la recapitulación, generada de terminos.json (PLAN-LIBRO.md §2.7).
\chapter{Glosario}
\label{ap:glosario}

Recapitulación de los términos técnicos, en el orden en que el libro los introduce. Cada entrada
lleva la etiqueta de \textbf{de qué lado está lo que nombra}, que es la primera pregunta que hay que
saber contestar de cualquier palabra de este libro.

\begin{center}\footnotesize
\nivelterm{meta} del lado de Lean, donde razonamos\quad
\nivelterm{objeto} del lado del sistema aritmético del que se habla\\[0.35em]
\nivelterm{meta→objeto} una función de Lean que \emph{produce} sintaxis de ese sistema\quad
\nivelterm{ambos} nombra la distinción misma\\[0.35em]
\nivelterm{proyecto} una noción sobre el proyecto, no sobre las teorías
\end{center}

\input{extraido/_glosario.tex}

%% Apéndice: notación. Recapitulación; los signos se introducen en la Parte I.
\chapter{Notación}
\label{ap:notacion}

\section{Los signos}


\begin{center}\small
\begin{tabularx}{\linewidth}{@{}l >{\raggedright\arraybackslash}X@{}}
\toprule
signo & qué es \\
\midrule
$\defnot{\Gamma \vdash A}$ & $A$ se deriva de las hipótesis $\Gamma$ en el cálculo $\omega$ \\
$\defnot{\Prf\,A}$ & $A$ se deriva en el cálculo de Hilbert finitario, sin hipótesis \\
$\sigma n$ & el sucesor de $n$ \\
$t \doteq u$ & igualdad \textbf{del lenguaje objeto} (no la de Lean, que es \ident{=}) \\
\texttt{\#0} & la variable ligada más cercana: un \defterm{índice de De Bruijn} \\
$\defnot{\codigo{\varphi}}$ & el \defterm{código de Gödel} de $\varphi$: un término objeto que la denota \\
$\defnot{\dotado{a}}$ & «$a$ dotado»: el código de $a$ construido \emph{dentro} del lenguaje (\ident{tcFn}) \\
$\defnot{\Prov(\codigo{\varphi})}$ & la fórmula objeto que dice «$\varphi$ es demostrable» \\
$\Delta_0$, $\Sigma_1$ & clases de fórmulas por su cuantificación: acotada, y $\exists$ sobre acotada \\
\bottomrule
\end{tabularx}
\end{center}

\nota{Ningún término técnico aparece en este libro antes de estar definido aquí. Lo comprueba
\modulo{doc/book/scripts/terminos.py} en cada compilación, y no admite excepciones.}

\section{Cómo se lee un nombre}

Los identificadores van en inglés y siguen las convenciones de Mathlib. Tres reglas bastan para
leer casi cualquiera:

\begin{enumerate}[leftmargin=1.4em]
\item \textbf{La conclusión primero, las hipótesis con \ident{_of_}.} El nombre dice \emph{qué se
  demuestra}, no cómo. \ident{isNat_succ_of_isNat} es «$\sigma n$ es natural, \emph{a partir de}
  que $n$ lo es».
\item \textbf{Los bicondicionales acaban en \ident{_iff}.} \ident{prf_In_iff_boundedIn} es un «si y
  sólo si».
\item \textbf{Los símbolos se deletrean}: \ident{add}, \ident{mul}, \ident{succ}, \ident{eq},
  \ident{lt}, \ident{le}, \ident{dvd}, \ident{not}, \ident{forall}, \ident{exists}. Así
  \ident{prf_add_zero_left} se lee sin traducir.
\end{enumerate}

Y hay algo más útil todavía, porque conecta con la chapa de niveles:
\textbf{el prefijo dice en qué nivel se afirma.} No es una impresión — es la práctica medida del
proyecto:

\begin{center}\small
\begin{tabularx}{\linewidth}{@{}l r >{\raggedright\arraybackslash}X@{}}
\toprule
prefijo & símbolos & qué anuncia \\
\midrule
\ident{ax_}   & 154 & un \defterm{axioma objeto}: se postula, no se demuestra \\
\ident{prf_}  & 585 & un teorema del cálculo finitario \ident{Prf} \\
\ident{pcc_}  & 214 & un teorema de \ident{Prf} \textbf{dentro de} \ident{Prov} — la teoría hablando de sí misma \\
\ident{CRIT_} & 13  & un \defterm{sondeo} crítico: una pregunta contestada, a menudo en negativo \\
\bottomrule
\end{tabularx}
\end{center}

\nota{Los recuentos son de \modulo{ROBINSON\_PlusPlus/} a 2026-09-03. La excepción conocida son las
seis meta-reglas $\omega$ de \modulo{FOL}, anteriores a la convención, que no llevan \ident{ax_}.}

\nota{Ningún término ni notación aparece en este libro antes de haberse introducido. Lo comprueba
\modulo{doc/book/scripts/terminos.py} en cada compilación, y no admite excepciones.}

\chapter{Licencia}
\label{ap:licencia}

Esta obra se distribuye bajo \textbf{dos} licencias, porque contiene dos clases de material.

\section{La prosa, las figuras y la maquetación}

\noindent\copyright{} 2026 Julián Calderón Almendros.

\medskip
\noindent Bajo \textbf{Creative Commons Reconocimiento-CompartirIgual 4.0 Internacional}
(\textbf{CC BY-SA 4.0}).

\begin{center}\small
\url{https://creativecommons.org/licenses/by-sa/4.0/deed.es}\\
Articulado completo: \url{https://creativecommons.org/licenses/by-sa/4.0/legalcode.es}
\end{center}

\noindent Resumen sin valor legal —manda el articulado—: cualquiera puede \textbf{copiar,
redistribuir, traducir y adaptar} esta obra, \textbf{incluso con fines comerciales}, siempre que
(a) \textbf{reconozca la autoría} y enlace a la licencia, indicando si hizo cambios, y (b) distribuya
sus obras derivadas \textbf{bajo esta misma licencia}.

\section{Los fragmentos de código Lean 4}

Los bloques de código de este libro se extraen de \modulo{ROBINSON\_PlusPlus/}, de
\modulo{sondeos/} y del proyecto \modulo{FOL}. \textbf{No cambian de licencia al ser citados}:
siguen bajo la licencia MIT del repositorio, que exige reproducir su aviso al distribuir porciones
sustanciales. Se reproduce aquí íntegro.

\begin{lean}
MIT License

Copyright (c) 2026 Julián Calderón Almendros

Permission is hereby granted, free of charge, to any person obtaining a copy
of this software and associated documentation files (the "Software"), to deal
in the Software without restriction, including without limitation the rights
to use, copy, modify, merge, publish, distribute, sublicense, and/or sell
copies of the Software, and to permit persons to whom the Software is
furnished to do so, subject to the following conditions:

The above copyright notice and this permission notice shall be included in all
copies or substantial portions of the Software.

THE SOFTWARE IS PROVIDED "AS IS", WITHOUT WARRANTY OF ANY KIND, EXPRESS OR
IMPLIED, INCLUDING BUT NOT LIMITED TO THE WARRANTIES OF MERCHANTABILITY,
FITNESS FOR A PARTICULAR PURPOSE AND NONINFRINGEMENT. IN NO EVENT SHALL THE
AUTHORS OR COPYRIGHT HOLDERS BE LIABLE FOR ANY CLAIM, DAMAGES OR OTHER
LIABILITY, WHETHER IN AN ACTION OF CONTRACT, TORT OR OTHERWISE, ARISING FROM,
OUT OF OR IN CONNECTION WITH THE SOFTWARE OR THE USE OR OTHER DEALINGS IN THE
SOFTWARE.
\end{lean}

\begin{observacion}
La doble licencia no es un tecnicismo: MIT es una licencia de \emph{software} y no dice nada sobre
traducciones, adaptaciones ni ediciones impresas, que es justo lo que una obra escrita necesita
regular. Y en sentido contrario, someter el código citado a CC BY-SA lo separaría de la licencia
bajo la que vive en el repositorio.
\end{observacion}


\backmatter
% \bibliographystyle{plain}
% \bibliography{bib/libro}

\end{document}
